\documentclass[11pt]{article}

\usepackage[T1]{fontenc}
\usepackage{lmodern}
\usepackage[a4paper,margin=27mm]{geometry}
\usepackage{amsmath,amssymb,amsthm,mathtools}
\usepackage{microtype}
\usepackage{enumitem}
\usepackage{booktabs,array}
\usepackage{xcolor}
\usepackage[colorlinks=true,linkcolor=blue!48!black,
  citecolor=blue!48!black,urlcolor=blue!55!black]{hyperref}
\setlength{\emergencystretch}{2em}

\newtheorem{theorem}{Theorem}[section]
\newtheorem{proposition}[theorem]{Proposition}
\newtheorem{corollary}[theorem]{Corollary}
\newtheorem{lemma}[theorem]{Lemma}
\theoremstyle{definition}
\newtheorem{definition}[theorem]{Definition}
\theoremstyle{remark}
\newtheorem{remark}[theorem]{Remark}

\newcommand{\C}{\mathbb C}
\newcommand{\R}{\mathbb R}
\newcommand{\N}{\mathbb N}
\newcommand{\cB}{\mathcal B}
\newcommand{\cD}{\mathcal D}
\newcommand{\cE}{\mathcal E}
\newcommand{\cG}{\mathcal G}
\newcommand{\cH}{\mathcal H}
\newcommand{\cM}{\mathcal M}
\newcommand{\cR}{\mathcal R}
\newcommand{\cS}{\mathcal S}
\newcommand{\Rea}{\operatorname{Re}}
\newcommand{\im}{\operatorname{im}}
\newcommand{\mul}{\operatorname{mul}}
\newcommand{\dom}{\operatorname{dom}}
\newcommand{\Div}{\operatorname{Div}}
\newcommand{\Spec}{\operatorname{Spec}}

\title{\textbf{The descent-before-polarization law for the Riemann divisor}\\
\large Source exactness, singular stability, and the forced change of carrier}
\author{Michael Haufschild}
\date{}

\begin{document}

\maketitle

\begin{abstract}
We identify a source-pointed exact-descent law behind the apparently
different source, cohomological, positive, and spectral approaches to the
Riemann divisor.  For an observation complex
\(E\xrightarrow O F\xrightarrow\delta G\), an actual source sector
\(S\subset E\), and an admissible sector \(A\subset F\), its three gates are
\[
 S\cap\ker O=0,\qquad O(S)=A\cap\ker\delta,\qquad
 \|Ox\|\ge c\|x\|\quad(x\in S)
\]
for some \(c>0\).  They are respectively uniqueness, existence, and stable
recovery of source lifts.  Every admissible compatible datum has a unique
source lift exactly when the first two identities hold.  Serial composition
is conservative exactly edgewise on propagated source sectors; parallel
observation intersects marginal kernels; stagewise lower bounds multiply.
In Hilbert geometry the remaining positive gate is equivalent to a
contractive factorization by Gram order.  Hence downstream positivity cannot
recreate information already killed upstream, and topology alone cannot
choose a Hermitian normalization.

For compositional conservation, finite augmentation has exactly the
universal zero-sum module as kernel.  Total covariance therefore controls
factor defects only on a source sector transverse to that kernel.

This law corrects a tempting Riemann target.  The finite arithmetic Green
operator factors through conjugate Riesz representation and odd Hodge
duality.  We prove that it is Frobenius-covariant unconditionally: its two
internal covariance defects are nonzero in general but cancel exactly as
\((u,-u)\).  At a divisor copy \(\rho\), one coordinate is the signed
boundary fracture
\[
                       \kappa(\rho)=2\Rea\rho-1.
\]
Consequently RH is equivalent to vanishing of the \emph{resolved} Green
defect, not to covariance of the composite.  More generally, RH follows
from one precise source statement: the actual resolved Green range must lie
in a completed arithmetic sector on which codiagonal augmentation is
conservative.

The complete shell source is finitely injective but cofinally collapsing;
its faithful normal reader is necessarily singular.  A raw M\"obius--theta
prime head retains label, unsquared orientation, and Xi multiplicity, and
lands coefficient-faithfully in completed Hellinger--Gamma incidence whose
local trace detects zero displacement.  A literal M\"obius--Gaussian Neumann
series realizes the forced Poisson lift on its absolute-convergence chamber.
The full Euler source already has four-unit Fisher energy, whereas its
division-free pole-cancelled reporter has limiting energy eight.  Its full
damping orbit is total: preservation of every coefficient uniquely fixes
that reporter, so a four-unit same-carrier compression is impossible.
Exact normal-crossing face descent and a relative Sullivan resolution
supply algebraic gluing, but neither chooses the required changed carrier.
The unresolved edge is therefore one global arithmetic descent that is
simultaneously exact, stable after a necessarily nonlocal singular
completion, and correctly polarized.  At multiplicity \(m\), its raw local
lift has pole order exactly \(m+1\).  All asserted statements are
unconditional; that global descent, and hence RH, is not claimed.
\end{abstract}

\medskip
\noindent\textbf{Keywords.}
Riemann hypothesis; source-pointed exact descent; resolved defect;
augmentation; co-Poisson transform; Meyer cohomology; Cartan exponent;
singular completion; contractive factorization; spectral synthesis.

\medskip
\noindent\textbf{2020 Mathematics Subject Classification.}
Primary 11M26, 11M06; Secondary 46N50, 47A06, 58A14.

\section{Introduction}

Many approaches to the Riemann hypothesis construct a source expansion,
a reflection, a positive form, a cohomological quotient, a recurrence
law, or a spectral operator.  These constructions look unrelated because
they discharge different gates of one source-pointed descent problem.
Separation makes a source lift unique; exactness glues every compatible
observation; a lower modulus makes the lift stable under completion; Gram
order calibrates its polarization.  Serial and parallel networks propagate
these obligations but do not identify them.  Augmentation is the decisive
Riemann instance: several source-resolved channels are combined into one
invariant, scalar, cohomological, or operator-valued output.  Equal and
opposite defects can then disappear exactly.

This elementary loss is decisive for the completed zeta function.  The
finite divisor Green map is already Frobenius-covariant at every divisor
point, on or off the critical line.  Its covariance is not evidence for
RH: it is exact cancellation of the Riesz and Hodge factor defects.  The
common invariant is therefore not merely the divisor or the split part of
its character, but the source-resolved defect before augmentation.

Every one-parameter complex character has the polar decomposition
\[
 e^{tz}=e^{t\mu}e^{it\tau},\qquad z=\mu+i\tau.
\tag{1.1}\label{eq:character-polar}
\]
The compact coordinate \(\tau\) governs phase and recurrence.  The split
Cartan coordinate \(\mu\) governs two-sided conditioning and positive
admission.  Passing to projective phase can remove a common compact phase;
it cannot change \(\mu\).  Passing to an invariant positive geometry can
retain a nonzero vector only when \(\mu=0\).  This elementary distinction
is the organizing principle of the paper.

For \(s\in\C\), define
\[
                    \kappa(s)=2\left(\Rea s-\frac12\right).
\tag{1.2}\label{eq:kappa}
\]
Critical-line reflection \(s\mapsto1-\overline{s}\) sends
\(\kappa\) to \(-\kappa\).  The functional equation therefore supplies
reflected pairs with opposite real rates.  It does not make either rate
zero.

The results answer four questions.

\begin{enumerate}[leftmargin=2.2em]
\item \emph{What is the parent law?}
      Source-pointed polarized exact descent is unique gluing plus stable,
      correctly normalized recovery.  Full coefficient preservation and
      strict Hermitian recalibration do not commute: total tests force the
      reporter and make every coefficient-preserving endomorphism the
      identity.  Exact descent must therefore precede polarization.  Serial conservativity is edgewise on
      propagated sectors; parallel conservativity is absence of a common
      marginal radical; Hilbert positivity is contractive factorization.
      Hence information must be resolved before a lossy quotient, never
      reconstructed after it.
\item \emph{Where is the Riemann Cartan coordinate before a zero is selected?}
      It is the logarithmic modulus of a cofinal prime-shell character.
      Two prime clocks recover the entire complex spectral coordinate.
\item \emph{Do the apparently different consumers read the same invariant?}
      Yes.  Reflected shell entropy, the nearby-cycle outer charge, and
      the completed Xi residue give its sign; Borel and Sonin growth give
      its absolute value; co-Poisson holonomy gives its odd hyperbolic
      transform.  Differential Cartan comparison is its higher-rank form,
      while quantum recurrence measures the independent compact coordinate.
      These are projections or consumers of one resolved source defect;
      none is automatically conservative on its cancellation module.
\item \emph{What theorem remains?}
      Local physical Xi jets now land coefficient-faithfully in completed
      Hellinger--Gamma incidence, where an unsquared trace detects modular
      displacement exactly.  Normal-crossing faces glue, the forced lift has
      a literal Neumann source before continuation, and four-unit source
      capacity exists.  Missing object is one global source-derived
      comparison with actual Green defect that first descends to the correct
      compatibility quotient, then remains exact through the forced nonlocal
      singular completion, and only there supplies the contractive Hermitian
      calibration, or a source law forcing that trace to vanish.
      Its scalar
      equivalents are boundedness of the centered Gaussian evaluation orbit
      and continuity of actual evaluations for the half-density quotient
      seminorm.
\end{enumerate}

The last statement is narrower than global spectral synthesis or a
Hilbert--P\'olya construction.  Positive completion is already available;
its continuous Xi spectrum is exactly the critical-line part of the
divisor.  The unresolved direction is a source-derived arithmetic descent,
not a new ambient space or a downstream scalar inequality: actual resolved
Green support has not yet been glued into a completed conservative
arithmetic sector with the required polarization.  Entry of every Xi
functional into continuous zero-Cartan support is the scalar form of that
statement.  Even universal continuity is stronger than logically
necessary, since a bounded scalar orbit of one explicit Gaussian witness at
each divisor point already forces RH.

\section{The source-pointed polarized exact-descent law}

Let \(T:E\to F\) be linear and let \(S\subset E\) be a source sector.  Say
that \(T\) is \emph{conservative on \(S\)} when

\[
                         S\cap\ker T=\{0\}.
\tag{2.0}\label{eq:conservative-on}
\]

This is relative injectivity, not global invertibility.  It asks exactly
whether \(T\) retains distinctions that the actual source can present.

\begin{theorem}[Source-pointed observation network law]
\label{thm:chain-law}
Let
\[
 E_0\xrightarrow{T_0}E_1\xrightarrow{T_1}\cdots
 \xrightarrow{T_{m-1}}E_m,
\qquad S_{i+1}=T_iS_i,
\]
where \(S_0\subset E_0\) is a submodule.  Then
\[
 T_{m-1}\cdots T_0\text{ is conservative on }S_0
 \quad\Longleftrightarrow\quad
 T_i\text{ is conservative on }S_i\text{ for every }i.
\tag{2.0a}\label{eq:chain-law}
\]
For normed modules, if constants \(c_i>0\) satisfy
\[
             \|T_ix\|\ge c_i\|x\|\qquad(x\in S_i),
\]
then
\[
 \|T_{m-1}\cdots T_0x\|
       \ge\left(\prod_{i=0}^{m-1}c_i\right)\|x\|
       \qquad(x\in S_0).
\tag{2.0b}\label{eq:chain-modulus}
\]
Thus failure at one edge cannot be repaired by any later transport or
reader.  Stable finite recovery is supplied by stagewise positive lower
moduli; finite-stage injectivity alone gives no uniform cofinal recovery.

For two parallel observations \(f:E\to F\) and \(h:E\to H\), let
\((f,h):E\to F\oplus H\) retain both outputs.  Then
\[
 \ker(f,h)=\ker f\cap\ker h,
 \qquad
 (f,h)\text{ is conservative on }S
 \Longleftrightarrow S\cap\ker f\cap\ker h=\{0\}.
\tag{2.0c}\label{eq:parallel-law}
\]
Serial postprocessing can only enlarge an invisible radical; parallel
resolution before loss replaces marginal radicals by their intersection.
\end{theorem}

\begin{proof}
For two stages \(E\xrightarrow fF\xrightarrow gG\),
\[
                   \ker(gf)=f^{-1}(\ker g).
\]
If \(gf\) is conservative on \(S\), then \(f\) is conservative on \(S\),
and \(g\) is conservative on \(fS\).  Conversely, \(gf(x)=0\) first gives
\(f(x)=0\) by conservativity of \(g\) on \(fS\), then \(x=0\) by
conservativity of \(f\) on \(S\).  Induction proves
\eqref{eq:chain-law}.  Successive substitution of the lower estimates gives
\eqref{eq:chain-modulus}.
For the parallel statement, \((f(x),h(x))=0\) if and only if both
coordinates vanish.  Finally, \(\ker f\subset\ker(gf)\), so no serial
postprocessor can remove a radical already present.
\end{proof}

Conservativity answers whether two source packets can become
indistinguishable.  It does not answer whether compatible observed data come
from the source.  The missing half is exact descent.

\begin{theorem}[Source-pointed exact-descent law]
\label{thm:exact-descent}
Let
\[
                 E\xrightarrow{O}F\xrightarrow{\delta}G
\]
be linear maps, and let \(S\subset E\), \(A\subset F\) be submodules.  The
following are equivalent:
\begin{enumerate}[leftmargin=2.2em]
\item \(O\) is conservative on \(S\) and
\[
                         O(S)=A\cap\ker\delta;
\tag{2.0d}\label{eq:exact-descent}
\]
\item the delivered source range lies in \(A\cap\ker\delta\), and every
      \(y\in A\) with \(\delta y=0\) has a unique lift \(x\in S\) satisfying
      \(Ox=y\).
\end{enumerate}
If additionally
\[
                         \|Ox\|\ge c\|x\|\qquad(x\in S)
\tag{2.0e}\label{eq:stable-descent}
\]
for one \(c>0\), then every compatible datum has a unique lift obeying
\(c\|x\|\le\|y\|\).  We call these three requirements
\emph{separation}, \emph{gluing}, and \emph{stability}.
\end{theorem}

\begin{proof}
Assume \eqref{eq:exact-descent}.  Membership in the right-hand side gives a
source lift.  If \(Ox=Ox'\) with \(x,x'\in S\), then
\(x-x'\in S\cap\ker O\), hence \(x=x'\).  Conversely, unique lifting of the
zero datum forces \(S\cap\ker O=0\), and existence for every member of
\(A\cap\ker\delta\) gives the reverse inclusion in
\eqref{eq:exact-descent}; the forward inclusion is the complex condition.
The norm estimate follows by substituting \(Ox=y\) into
\eqref{eq:stable-descent}.
\end{proof}

In Hilbert spaces stability has a polarized form.  For bounded readers
\(A_0,B_0:H\to K\), Douglas factorization gives
\[
 B_0B_0^*-A_0A_0^*\ge0
 \quad\Longleftrightarrow\quad
 \exists C:H\to H,\qquad \|C\|\le1,\qquad A_0=B_0C.
\tag{2.0f}\label{eq:douglas-descent}
\]
Thus positive Gram order is precisely a contractive descent factor, not a
separate scalar miracle.  Yet topology and Hermitian calibration remain
independent: the topology induced by a nonzero conormal \(L\) is unchanged
under \(L\mapsto aL\), while \(\|aLx\|^2=|a|^2\|Lx\|^2\).  Projective source
topology can locate the conormal line but cannot select a numerical budget
such as four units.

There is a further order constraint.  Let \(H\) be a complex inner-product
space and let \((q_i)_{i\in I}\subset H\) be total, meaning

\[
 \bigcap_{i\in I}\{z:\langle z,q_i\rangle=0\}=\{0\}.
\tag{2.0g}\label{eq:total-tests}
\]

\begin{theorem}[Descent before polarization]
\label{thm:descent-before-polarization}
For \(x,y\in H\), equality of the complete coefficient packets

\[
             \langle y,q_i\rangle=\langle x,q_i\rangle
             \qquad(i\in I)
\tag{2.0h}\label{eq:coefficient-packet}
\]

forces \(y=x\).  Consequently, if \(B<\|x\|^2\), no same-carrier reporter
can preserve \eqref{eq:coefficient-packet} and have squared norm at most
\(B\).  More strongly, every linear endomorphism of \(H\) preserving all
coefficients of every vector is the identity.
\end{theorem}

\begin{proof}
The vector \(y-x\) lies in every left-orthogonal hyperplane in
\eqref{eq:total-tests}, hence is zero.  The norm obstruction follows by
substitution.  Applying the first assertion to \(Cx\) and \(x\), for every
\(x\), proves \(C=1\).
\end{proof}

Thus exactness and polarization do not generally commute.  If a total
pre-descent packet has the wrong energy, no choice of another vector and no
nontrivial preprocessing endomorphism of the same carrier can correct it
while retaining every coefficient.  One must first change carrier or
incidence by an exact quotient which discards only compatibility-redundant
data; the Hermitian Gram order is then posed on that descended carrier.
This is the \emph{descent-before-polarization law}.

Every route in this paper is therefore audited by one operation: propagate
its actual source sector, test separation, test gluing of compatible data,
and then test stable and Hermitian recovery.  Retain necessary parallel
coordinates before loss; no later reader can repair an earlier failure.
Resolved augmentation is the central Riemann specialization.  Cartan
growth, positive admission, and vertical completion are rank-one,
polarized, and topological shadows of the same law.

\subsection{Rank-one Cartan consequence}

Let \(E\) be a complex vector space, \(V_t:E\to E\) a family of
complex-linear maps, and suppose one nonzero source direction satisfies
\[
                         V_tx=e^{tz}x,
             \qquad z=\mu+i\tau.
\tag{2.1}\label{eq:vector-character}
\]
No group law, topology, completeness, or spectral resolution is needed.

\begin{theorem}[Pointed Cartan conservation and fate]
\label{thm:cartan-fate}
For every complex seminorm \(p\) and every \(t\in\R\),
\[
                         p(V_tx)=e^{t\mu}p(x).
\tag{2.2}\label{eq:primal-growth}
\]
Consequently, for \(t\ge0\),
\[
 \max\{p(V_tx),p(V_{-t}x)\}=e^{t|\mu|}p(x).
\tag{2.3}\label{eq:reflected-condition}
\]
If \(p\) is invariant on the orbit, then
\[
                         p(x)=0\quad\hbox{or}\quad\mu=0.
\tag{2.4}\label{eq:primal-alternative}
\]

Dually, let \(L\ne0\) satisfy
\[
                         L(V_tu)=e^{tz}L(u).
\tag{2.5}\label{eq:dual-character}
\]
There is a finite nonnegative invariant gauge \(g\) and \(C\ge0\) with
\(|L(u)|\le Cg(u)\) for every \(u\) if and only if \(\mu=0\).

Suppose now that \(p(V_tu)=p(u)\) for all \(t,u\), and choose any \(u\)
with \(L(u)\ne0\).  The rescaled transports
\[
 u_n=e^{-t_nz}V_{t_n}u,\qquad t_n=n\mu,
\tag{2.6}\label{eq:graph-null-sequence}
\]
obey the exact laws
\[
 L(u_n)=L(u),\qquad
 p(u_n)=e^{-n\mu^2}p(u).
\tag{2.7}\label{eq:graph-null-law}
\]
If \(\mu\ne0\), the graph of \(L\) therefore acquires the nonzero vertical
limit \((0,L(u))\) in the \(p\)-completion.  For scalar output its closed
linear graph relation contains \(\{0\}\times\C\): it is maximally vertical
over zero and cannot be single-valued.

Finally, on the real reciprocal plane
\[
                         (a,b)\longmapsto
                         (e^{t\mu}a,e^{-t\mu}b),
\tag{2.8}\label{eq:reciprocal-plane}
\]
every invariant symmetric quadratic form is a multiple of \(ab\) when
\(\mu\ne0\), and every invariant positive-semidefinite quadratic form is
zero.
\end{theorem}

\begin{proof}
Absolute homogeneity of \(p\), together with
\(|e^{tz}|=e^{t\Rea z}\), proves \eqref{eq:primal-growth}.  Apply it at
\(t\) and \(-t\), then use
\(\max(e^{t\mu},e^{-t\mu})=e^{t|\mu|}\), to obtain
\eqref{eq:reflected-condition}.  If \(p(x)\ne0\) and the orbit is
invariant, \eqref{eq:primal-growth} at \(t=1\) gives \(e^\mu=1\), hence
\(\mu=0\).

For the dual statement, choose \(u\) with \(L(u)\ne0\).  A dominating
invariant gauge would bound \(e^{t\mu}|L(u)|\) in both time directions,
which is possible only at \(\mu=0\).  Conversely, if \(\mu=0\), take
\(g(u)=|L(u)|\).  In \eqref{eq:graph-null-sequence}, linearity and the
dual character law cancel the two exponential factors, while absolute
homogeneity and invariance give
\(p(u_n)=e^{-t_n\mu}p(u)=e^{-n\mu^2}p(u)\).  This tends to zero for
\(\mu\ne0\).  Scaling the nonzero vertical limit produces every scalar.
For the reciprocal plane, invariance of
\(Aa^2+2Bab+Db^2\) forces \(A=D=0\); positivity at \((1,1)\) and
\((1,-1)\) then forces \(B=0\).
\end{proof}

Equation~\eqref{eq:primal-alternative} is the decisive direction: an
equivariant comparison into invariant positive geometry cannot rotate a
nonzero split character into a unitary one.  It either preserves zero
Cartan slope or sends the source direction to the target radical.  We call
a nonzero dual character \emph{temperately admitted} when the gauge in the
theorem exists.  Thus tempered admission is precisely zero-Cartan
admission, not a new spectral construction.

\subsection{Algebraic descent}

Let \(X,Y,Z\) be modules, let \(C,B,A\) be endomorphisms, and let
\(X\xrightarrow rY\xrightarrow jZ\) be linear maps.  Define the factor
defects
\[
       \partial r=Br-rC,\qquad \partial j=Aj-jB
\tag{2.9}\label{eq:factor-defects}
\]
and the resolved composite defect
\[
 \Delta_{j,r}(x)=
 \left((\partial j)(rx),\,j((\partial r)x)\right)\in Z\oplus Z.
\tag{2.10}\label{eq:resolved}
\]
Write \(\Sigma(u,v)=u+v\).

\begin{proposition}[Resolved-defect conservation]
\label{prop:resolved}
For every \(x\in X\),
\[
                  A(jr)x-(jr)Cx=\Sigma\Delta_{j,r}(x).
\tag{2.11}\label{eq:defect-conservation}
\]
Moreover,
\[
 0\longrightarrow Z\xrightarrow{z\mapsto(z,-z)}Z\oplus Z
 \xrightarrow{\Sigma} Z\longrightarrow0
\tag{2.12}\label{eq:augmentation}
\]
is exact.  If the two coordinates of \(\Delta_{j,r}\) lie in submodules
\(P,Q\subset Z\) with \(P\cap Q=0\), then covariance of the composite
forces both resolved coordinates to vanish.
\end{proposition}

\begin{proof}
Insert and subtract \(jBrx\):
\[
 Ajrx-jrCx=(Aj-jB)rx+j(Br-rC)x.
\]
The kernel of \(\Sigma\) consists exactly of the pairs \((z,-z)\).
If \(u\in P\), \(v\in Q\), and \(u+v=0\), then \(u=-v\in P\cap Q\);
hence \(u=v=0\).
\end{proof}

The two-channel statement is the first case of an exact finite law.  For
\(n\ge1\), set
\[
 \Sigma_n:Z^n\longrightarrow Z,\qquad
 \Sigma_n(z_0,\ldots,z_{n-1})=\sum_{i=0}^{n-1}z_i,
\tag{2.12a}\label{eq:finite-augmentation}
\]
and define
\[
 \partial_n(w_0,\ldots,w_{n-2})
   =\left(-\sum_{i=0}^{n-2}w_i,w_0,\ldots,w_{n-2}\right).
\tag{2.12b}\label{eq:finite-boundary}
\]

\begin{theorem}[Resolved source-conservativity]
\label{thm:augmentation-conservativity}
For every module \(Z\),
\[
0\longrightarrow Z^{n-1}\xrightarrow{\partial_n}Z^n
 \xrightarrow{\Sigma_n}Z\longrightarrow0
\tag{2.12c}\label{eq:finite-exact}
\]
is exact.  Hence
\[
                 \ker\Sigma_n=\im\partial_n.
\tag{2.12d}\label{eq:finite-kernel}
\]
If \(S\subset Z^n\) is a submodule, the following are equivalent:
\[
 S\cap\ker\Sigma_n=\{0\};
 \qquad
 \text{for every }D:X\to Z^n\text{ with }\im D\subset S,
 \quad \Sigma_nD=0\Longleftrightarrow D=0.
\tag{2.12e}\label{eq:source-conservative}
\]
No observation applied after \(\Sigma_n\) can
distinguish a nonzero vector in \(S\cap\ker\Sigma_n\).

For endomorphisms \(C,T_0,\ldots,T_{n-1}\) of one module, put
\(\delta_C(T)=CT-TC\).  The defect of their ordered product is
\[
 \delta_C(T_{n-1}\cdots T_0)
 =\sum_{i=0}^{n-1}
  T_{n-1}\cdots T_{i+1}\,\delta_C(T_i)\,
  T_{i-1}\cdots T_0,
\tag{2.12f}\label{eq:finite-leibniz}
\]
with empty products equal to the identity.  Thus
\eqref{eq:source-conservative} is exactly the condition under which
covariance of a product controls its transported factor defects.
\end{theorem}

\begin{proof}
The last \(n-1\) coordinates of \(\partial_n w\) recover \(w\), so
\(\partial_n\) is injective, and its coordinate sum is zero.  Conversely,
if \(\sum_i z_i=0\), take \(w_i=z_{i+1}\); then
\(z_0=-\sum_iw_i\), proving \eqref{eq:finite-kernel}.  The vector
\((z,0,\ldots,0)\) proves surjectivity of \(\Sigma_n\).

Now \(\Sigma_nD=0\) places every value of \(D\) in
\(S\cap\ker\Sigma_n\).  This forces \(D=0\) exactly when that intersection
is zero.  A post-augmentation observer still sends every element of
\(\ker\Sigma_n\) to zero.  Finally, expand the right side of
\eqref{eq:finite-leibniz}; all intermediate terms cancel in adjacent pairs,
leaving \(C(T_{n-1}\cdots T_0)-(T_{n-1}\cdots T_0)C\).
\end{proof}

Thus composite covariance is a quotient statement.  It may hold because
several nonzero internal defects form a zero-sum packet.  A source solves
the inference problem only by placing the actual packet in a sector where
that cancellation is impossible.

The same law has a genuine gluing realization at a normal crossing.  Let
\(P=R[X_i:i\in I]\) for a finite set \(I\), put
\(I_i=(X_i)\), \(I_{ij}=(X_i,X_j)\), and
\(J=(\prod_{i\in I}X_i)\).  Then
\[
 J=\bigcap_{i\in I}I_i,
\tag{2.12g}\label{eq:face-uniqueness}
\]
and a family \((a_i\bmod I_i)_i\) comes from one class in \(P/J\) whenever
\[
                         a_i-a_j\in I_{ij}
                         \qquad(i,j\in I).
\tag{2.12h}\label{eq:face-compatibility}
\]
Thus atomic-face restriction is injective and every pairwise compatible
family glues.  This is Theorem~\ref{thm:exact-descent} with both algebraic
gates discharged, not merely a family of local reconstructions.  The
squarefree product is the complete ambiguity ideal, while pairwise faces
are the compatibility reader.

Relative Sullivan attachment is the homological version of the same law.
Over a characteristic-zero field, start from the free degree-one
graded-commutative algebra on a basis of \(V\), augmented onto the
square-zero algebra \(k\oplus V\).  Each successor
stage attaches generators precisely to cocycles killed by the augmentation.
The limit augmentation is surjective and a quasi-isomorphism.  It therefore
constructs gluing by killing the compatibility obstruction, whereas an
unframed resolution still need not separate source presentations that have
the same target class.  Exactness and source conservativity are independent
even in a universal resolution.

\subsection{Closed relations and verticalization}

For a linear relation \(\cR\subset X\times Y\), define
\[
 \mul\cR=\{y:(0,y)\in\cR\},\qquad
 \dom\cR=\{x:\exists y,\ (x,y)\in\cR\}.
\]

\begin{proposition}[Closed retention dichotomy]
\label{prop:closed-retention}
Let \(X,Y\) be Hausdorff topological vector spaces and let
\(\cR\subset X\times Y\) be a closed linear relation.

\begin{enumerate}[label=\textup{(\roman*)},leftmargin=2em]
\item If \((x_n,y_n)\in\cR\), \(x_n\to0\), and \(y_n\to y\), then
      \(y\in\mul\cR\).
\item If the closed linear span of all such output limits is \(Y\), then
      \(\mul\cR=Y\) and
      \[
                        \cR=\dom\cR\times Y.
      \tag{2.13}\label{eq:vertical}
      \]
\end{enumerate}
\end{proposition}

\begin{proof}
Closedness sends the limit of \((x_n,y_n)\) to
\((0,y)\in\cR\).  Therefore every output limit belongs to
\(\mul\cR\).  This multivalued part is closed and linear, so a dense
linear family inside it implies \(\mul\cR=Y\).  If
\((x,y_0)\in\cR\), then \((0,y-y_0)\in\cR\) for every \(y\in Y\);
linearity gives \((x,y)\in\cR\), which proves
\eqref{eq:vertical}.
\end{proof}

The theorem gives the topological analogue of
\eqref{eq:augmentation}.  A regular single-valued map sends a collapsing
input difference to zero.  A closed relation which retains densely many
limits avoids erasure only by accepting every output at zero; it then
ceases to transport output information.

\subsection{Higher-rank conservation}

The rank-one law is the one-dimensional shadow of a prefix theorem.  Let
\(\sigma_{n,k}\) and \(\tau_{n,k}\) be integer-valued prefix invariants of
two finite-rank systems, normalized at \(k=0\).  If one regular comparison
gives a constant \(C\), independent of \(n,k\), such that
\[
                         |\sigma_{n,k}-\tau_{n,k}|\le C,
\tag{2.14}\label{eq:bounded-prefix}
\]
then every finite normalized Cartan vector formed from prefix increments
has the same limit in the two systems:
\[
 \frac{\sigma_{n,k+1}-\sigma_{n,k}}n\longrightarrow L_k
 \quad\Longrightarrow\quad
 \frac{\tau_{n,k+1}-\tau_{n,k}}n\longrightarrow L_k.
\tag{2.15}\label{eq:cartan-prefix-limit}
\]
We call such a comparison \emph{Cartan-regular}.  This is the exact scope
of the term in this paper; no assertion is made about comparisons without
\eqref{eq:bounded-prefix} or exact rank-one equivariance.
Indeed, telescoping identifies the first \(k\) increment sums with the
corresponding prefixes, and the difference divided by \(n\) tends to zero.

This hypothesis is not formal decoration in the hypertranscendental
setting.  A finite matrix of differential operators with independent rows
has a right inverse over the skew field of rational pseudodifferential
operators.  Its row lattice has finite conductor in its saturation.
Intrinsic and ambient relative-index prefixes therefore differ by at most
twice this fixed conductor dimension, and
\eqref{eq:cartan-prefix-limit} applies.  Regular finite differential
re-embedding cannot compress or manufacture a normalized Cartan
direction.

\subsection{The universal fate law}

Theorem~\ref{thm:chain-law}, Theorem~\ref{thm:cartan-fate}, and
Theorem~\ref{thm:augmentation-conservativity}, together with
Propositions~\ref{prop:resolved}--\ref{prop:closed-retention}, give one
law:

\begin{quote}
\emph{Information survives a source-pointed network exactly when serial
edges are conservative on delivered sectors and parallel readers have no
common source radical.  Parallel resolution can protect a distinction
before loss; serial postprocessing cannot recreate it afterward.
Resolved conservation fails at its augmentation edge precisely when the
actual factor packet meets the zero-sum kernel.  Completion replaces exact
kernel intersection by a uniform lower-modulus problem.  Once an edge
kills a distinction, no later invariant, positive form, quotient, or reader
can recreate it; faithful cofinal recovery must instead become singular.}
\end{quote}

The parent assertion is the network law
\eqref{eq:chain-law}--\eqref{eq:parallel-law};
\eqref{eq:source-conservative} is its augmentation edge.  The Cartan fates
classify realized algebraic or stable-conservativity failures.  Single-valued positive descent is governed by
\eqref{eq:primal-alternative}; two-sided observation by
\eqref{eq:reflected-condition}; reciprocal bilinear retention by
\eqref{eq:reciprocal-plane}; algebraic factorization by the resolved defect;
and closed relational retention by Proposition~\ref{prop:closed-retention}.
The remaining sections identify the source sectors and augmentation maps
for the Riemann and hypertranscendental constructions.  Algebraic Xi
spectrum is abundant; positive continuous Xi spectrum is exactly the
zero-Cartan sector, but that fact becomes RH-producing only after a
source-support inclusion into an augmentation-conservative completed
sector.

\section{The complete source and its normal cone}

\subsection{Longitudinal collapse and projective survival}

Let \(S_K(s)\in\ell^1(\N,\C)\) denote the complete high-exponent source
beginning at shell depth \(K\), defined for \(\Rea s>0\).  The exact
longitudinal statement is
\[
                     S_K(s)\longrightarrow e_0
                     \qquad(K\to\infty)
\tag{3.1}\label{eq:source-collapse}
\]
for every fixed \(s\).  Hence any two fixed spectral points become
indistinguishable in the ordinary source norm.

The first two normalized adjacent-shell ratios define a projective chart
\[
 \mathfrak B_K(s)=
 \left(\mathfrak b_{2,K}(s),\mathfrak b_{3,K}(s)\right).
\tag{3.2}\label{eq:blowup-chart}
\]
Put
\[
 \chi_p(s)=p^{1/2-s},\qquad
 \mathfrak B_\infty(s)=\left(\chi_2(s),\chi_3(s)\right).
\tag{3.3}\label{eq:two-prime}
\]

\begin{proposition}[Collapse--blow-up duality]
\label{prop:blowup}
For every \(s\) with \(\Rea s>0\),
\[
 S_K(s)\to e_0,\qquad
 \mathfrak B_K(s)\to\mathfrak B_\infty(s).
\tag{3.4}\label{eq:collapse-blowup}
\]
The map \(\mathfrak B_\infty:\C\to\C^2\) is injective.  No uniformly
Lipschitz observer of the ordinary \(\ell^1\) source can retain a fixed
nonzero separation between two distinct limiting charts.
\end{proposition}

\begin{proof}
The head coefficient tends to one and the summable shell tail tends to
zero, giving \eqref{eq:source-collapse}.  Adjacent shell asymptotics give
\(\mathfrak b_{2,K}(s)\to2^{1/2-s}\) and
\(\mathfrak b_{3,K}(s)\to3^{1/2-s}\).

Equality of the first coordinates determines \(\Rea s\) and determines
\(\Im s\) modulo \(2\pi/\log2\).  Equality of the second coordinates
gives the analogous congruence modulo \(2\pi/\log3\).  Since
\(\log3/\log2\) is irrational, the two congruences have only the common
zero period.  Thus \(s\) is determined.

Finally, if \(T_K\) have a common Lipschitz constant \(L\), then
\[
 \|T_KS_K(s)-T_KS_K(t)\|
 \le L\|S_K(s)-S_K(t)\|\longrightarrow0.
\]
Any fixed separation of the limiting projective charts therefore forces
inverse conditioning to diverge.
\end{proof}

This proves that the spectral coordinate is not absent.  It lives in the
normal cone of the collapsing source.

\subsection{The exact finite-fibre boundary}

Let \(p_i\) be the increasing sequence of primes and put
\[
             P_j=(p_i^{-(j+1)})_{i\ge0}\in\ell^2(\N,\C).
\tag{3.5}\label{eq:prime-packet}
\]
Then \(P_j\to0\), but the packets are maximally cyclic.

\begin{theorem}[Finite-fibre strictness]
\label{thm:finite-fibre}
Every infinite subsequence of \((P_j)\) has dense linear span in
\(\ell^2\).  The same is true for every affine family
\(\{P_j-y:j\in J\}\) with \(J\subset\N\) infinite.

Let \(g:\N\to G\) assign an exponent to a grade and let \(y_\gamma\in
\ell^2\) be arbitrary centers.  A closed gradewise-proper carrier in
\[
                  \prod_{\gamma\in G}\ell^2
\]
which retains every supported packet
\(\mathbf 1_{g(j)}(P_j-y_{g(j)})\) exists if and only if every fibre
\(\{j:g(j)=\gamma\}\) is finite.
\end{theorem}

\begin{proof}
Fix \(k\).  Multiply \(P_j\) by \(p_k^{j+1}\) and subtract the previously
recovered coordinates.  The result is \(e_k\) plus a tail whose norm is
bounded by a geometric sequence with ratio
\(p_k/p_{k+1}<1\).  Along any infinite \(J\), the exponents tend to
infinity, so the tail tends to zero.  Induction on \(k\) places every
coordinate vector in the closure of the restricted packet span.
Orthogonal-complement duality gives the affine version.

If a grade fibre is infinite, its affine packet span is dense.  Closedness
then forces the entire independently supported \(\ell^2\) slice into the
carrier, contradicting gradewise properness.  Conversely, finite fibres
span finite-dimensional closed proper subspaces in each grade; their
product is a closed gradewise-proper carrier containing all packets.
\end{proof}

The theorem rules out an ungraded Hilbert or product completion as the
missing carrier.  Local finiteness is not a convenient technical
assumption; it is the exact boundary between a proper graph and saturated
verticality.

\subsection{Source realizations and their first lossy edge}

The same transverse coordinate appears before any divisor is selected.

\paragraph{Hypertranscendental shell frame.}
The high-exponent shell multipliers are differentially independent over
the completed-zeta differential field, have disjoint coefficient
channels, and are zero-free in the strict right half-plane.  Multiplying
the completed zeta function by any sufficiently deep shell replicates its
divisor with multiplicity.  Thus divisor data are longitudinal, while
\(\mathfrak B_\infty(s)\) is the labeled transverse coordinate.  Scalar
erasure of the label is nonclosable.

\paragraph{M\"obius--theta physical jets.}
For every \(k\ge0\) there is a Schwartz source \(f_k\), constructed by
internal polynomial promotion of one theta source, with
\[
                         \widehat f_k(t)=\Xi^{(k)}(t)
\tag{3.6}\label{eq:theta-landing}
\]
on the real spectral axis.  At a zero \(\rho\) of multiplicity \(m\),
the \(m\times m\) local moment matrix is triangular with
\(\Xi^{(m)}(\rho)\ne0\) on every diagonal entry.  Hence it is nonsingular.
No local multiplicity direction is absent from the source.

The result is now cofinal on the complex landing domain, not merely a
finite local span.  The order-\(k\), exponent-\(M\) source transform has the
form
\[
 \mathcal M_{k,M}(w)
   =\zeta\!\left(Ms(w)\right)^{-1}\Xi^{(k)}(w),
 \qquad
 \mathcal M_{k,K+N}(w)\longrightarrow\Xi^{(k)}(w),
\tag{3.6a}\label{eq:theta-cofinal-landing}
\]
throughout the honest theta strip with \(\Rea s(w)>0\).  The adjacent
division-free conormal pair lands simultaneously:
\[
 \bigl(\mathcal M_{k,K+N}(w),\mathcal M_{k+1,K+N}(w)\bigr)
 \longrightarrow\bigl(\Xi^{(k)}(w),\Xi^{(k+1)}(w)\bigr).
\tag{3.6b}\label{eq:theta-conormal-landing}
\]
At a zero \(\rho\) of multiplicity \(m\), its leading pair is the nonzero
vertical conormal
\[
                 \bigl(0,\Xi^{(m)}(\rho)\bigr).
\tag{3.6c}\label{eq:theta-leading-conormal}
\]
The normalized source cofaces telescope to the same rungs, so their finite
codiagonals converge to every Xi jet at every actual divisor coordinate.
This completes target-free complex source landing with multiplicity.

More strongly, one raw prime head retains arithmetic label, reflection
orientation, and multiplicity in the same division-free graph.  For a prime
\(p\) and integer \(d\ge1\), its internally promoted transform is
\[
 \mathcal L R_{k,p,d}(w)=A_{p,d}(s(w))\,\Xi^{(k)}(w),
 \qquad A_{p,d}(s)\ne0,
\tag{3.6d}\label{eq:prime-raw-landing}
\]
and its unsquared reflected modulus ratio is
\[
 \frac{|A_{p,d}(1-s)|}{|A_{p,d}(s)|}
       =p^{d(2\Rea s-1)}.
\tag{3.6e}\label{eq:prime-raw-orientation}
\]
A singular source-normal reader removes the half-density and logarithmic
translation and recovers the oriented local theta jet exactly.  At a zero
\(\rho\) of multiplicity \(m_\rho\), the first conormal pair is
\[
 \left(0,A_{p,d}(\rho)\Xi^{(m_\rho)}(\rho)\right)\ne(0,0).
\tag{3.6f}\label{eq:prime-raw-conormal}
\]
Thus neither squaring nor projectivization is needed to put source label,
Cartan sign, and divisor multiplicity in one physical graph.

Nor is the raw head an isolated surrogate.  Its complete prime-normal
residual multiplier \(R_{p,d}(w)\), synthesized before evaluation from all
arithmetic channels \(n\ge p\), obeys
\[
                         R_{p,d}(w)\longrightarrow\mu(p)
                         \qquad(d\to\infty)
\tag{3.6fa}\label{eq:prime-isolation}
\]
throughout its right-half-plane landing domain.  For prime \(p\), the limit
is nonzero.  Hence the complete residual source, not merely its isolated
head, eventually gives the nonzero leading vertical conormal at every Xi
divisor.  Arithmetic nonannihilation is source-derived by prime isolation;
it is not an additional scalar hypothesis.

The next edge is also available locally.  Let
\(c_\rho=\Xi^{(m_\rho)}(\rho)\ne0\).  For every distinct prime pair there is
a completed Calder\'on--Witt incidence vector \(I_\rho\) whose normalized
incidence coefficient is exactly \(c_\rho\), whose decoded source value is
\(c_\rho\) times the actual Gamma witness, and which is recovered by the
completed adjacent Hellinger cofaces.  If \(H_{a,n}(I_\rho)\) denotes its
gauged unsquared pure-Hodge realization, then
\[
 \operatorname{Tr}_{\rm sec}H_{a,n}(I_\rho)
     =c_\rho\,\Gamma_{a,n},
 \qquad
 \operatorname{Tr}_{\rm sec}H_{a,n}(I_\rho)=0
     \Longleftrightarrow a=0.
\tag{3.6g}\label{eq:theta-hodge-detector}
\]
Thus physical multiplicity provenance and a conservative orientation
detector now meet coefficient-faithfully before completion.  What remains
is not this local incidence bridge.  It is a global, source-derived
comparison identifying its displacement with the actual resolved Green
coordinate and a completed law forcing the unsquared secondary trace to
vanish.  Raw coface recovery alone holds for every \(a\), on and off axis.

\paragraph{Euler valuation and Hellinger--Gamma source.}
For \(q=p^{-\sigma}\), the diagonal activation
\(\zeta(\sigma)/\zeta(2\sigma)\) has local valuation law
\[
 \nu_q(0)=\frac1{1+q},\qquad
 \nu_q(k)=\frac{(1-q)q^k}{1+q}\quad(k\ge1).
\tag{3.7}\label{eq:valuation}
\]
Radical occupancy is Bernoulli; conditional positive valuation is
geometric.  The occupancy Fisher variance is \(q/(1+q)^2\), so
Fisher whitening necessarily amplifies rare prime events.  The resulting
Gamma normalization is an inverse-conditioned normal coordinate, not a
missing arithmetic source.

There is nevertheless an exact four-unit source before completion.  Unique
factorization \(n=m^2d\), with \(d\) squarefree, identifies the full positive
Euler Gibbs law with the independent product of squarefree occupancy and
the square-residual law.  Its direct projective Fisher energy satisfies
\[
       \mathcal I_{\rm Euler}(\sigma)\longrightarrow4
       \qquad(\sigma\downarrow1).
\tag{3.7aa}\label{eq:euler-four-unit}
\]
The physical affine Euler clock and the essential CCM cusp boundary layer
both converge to the unit exponential clock and carry the same four-unit
normalization; the complementary packet Ward flux is an asymptotic
isometry on that distinguished edge field.  Four-unit capacity is therefore
already source-derived and is not the missing datum.

What fails is the naive attempt to use the same carrier as a division-free
completed reporter.  If \(P_\sigma\) is the literal pole-cancelled
pre-scalar Gamma source, then
\[
 \|P_\sigma\|^2
   =4(\sigma-1)^2M_2(\sigma)\longrightarrow8.
\tag{3.7ab}\label{eq:prescalar-eight}
\]
For every damped Euler state its coefficient is exactly the un-divided
shifted prime numerator, and after the affine Gamma term is included the
completed numerator factors as
\[
 N_{\Xi}(\sigma,u)
   =G(\sigma,u)\,\mathcal T_{\Xi}
      \bigl(\sigma+u(\sigma-1)\bigr),
\tag{3.7ac}\label{eq:prescalar-factorization}
\]
where \(G\) is the retained vanishing graph denominator and
\(\mathcal T_{\Xi}\) is the continued prime--archimedean trace.  Thus the
division-free reporter exists, but it has an exact factor-two energy excess.
This excess is rigid.  The full nonnegative damping orbit
\((D_{\sigma,u})_{u\ge0}\) is total in the positive-Euler Hilbert carrier:
\[
 \bigl[\langle y,D_{\sigma,u}\rangle
       =\langle P_\sigma,D_{\sigma,u}\rangle\quad(u\ge0)\bigr]
       \Longrightarrow y=P_\sigma.
\tag{3.7ad}\label{eq:euler-reporter-rigidity}
\]
In particular, for all \(\sigma>1\) sufficiently near one, every reporter
preserving the full division-free coefficient family has squared norm
strictly greater than four.  Theorem~\ref{thm:descent-before-polarization}
therefore rules out a same-carrier four-unit compression.  Since the
conormal topology is projective, no topological refinement can select one
either.  The required map must first change the carrier or incidence by
exact descent, and only then satisfy a Gram domination and contractive
factorization as in \eqref{eq:douglas-descent}.  Which coefficient
directions become exact compatibility boundaries is the unresolved
arithmetic question.

At fixed primes, aperture, and depth, write the balanced completed
Hellinger--Gamma primitive as
\[
             \mathcal H_{a,n}(x)=(g_{-a,n}x,g_{a,n}x),
             \qquad g_{-a,n},g_{a,n}>0.
\tag{3.7a}\label{eq:balanced-hellinger}
\]
Its codiagonal is \((g_{-a,n}+g_{a,n})x\).  Hence
\[
 \im\mathcal H_{a,n}\cap\ker\Sigma_2=\{0\}.
\tag{3.7b}\label{eq:hellinger-transverse}
\]
Every finite balanced grade is therefore an
augmentation-conservative source sector.  In particular, a nonzero Green
antidiagonal cannot equal one finite balanced Hellinger--Gamma primitive.
This is a positive transversality theorem and a support no-go at once.

Restoring the grade label does not automatically solve the support problem.
Let \(v_n(x)\) place \(\mathcal H_{a,n}(x)\) in the \(n\)-th orthogonal
depth grade.  The Gamma-factorial gauge gives \(v_n(x)\to0\).  If a
secondary boundary \(B\) is defined on every \(v_n(x)\) and annihilates
every adjacent coface \(v_{n+1}(x)-v_n(x)\), then
\[
 Bv_n(x)=Bv_0(x)\quad\hbox{for every }n,
 \qquad (0,Bv_0(x))\in\overline{\operatorname{graph}B}.
\tag{3.7c}\label{eq:hellinger-vertical}
\]
Thus a nonzero endpoint charge makes the one-marginal graph nonclosed.
Finite transversality and completed support are separate obligations; the
successful sector must join factor provenance with the independent
physical source before this closure is formed.

\paragraph{Hypertranscendental character germs.}
The genuine exponential source clock has a canonical valuation--unit split
in the convergent local ring \(A=\C\{\lambda\}\):
\[
 e^{t\lambda}-1=\lambda U_t(\lambda),
 \qquad U_t(0)=t,
\tag{3.7d}\label{eq:clock-unit-split}
\]
and \(U_t\) is a unit whenever \(t\ne0\).  For a nonvanishing frame
\(\tau_0,\ldots,\tau_r\), the \(i\)-th physical coface generator is
therefore
\[
 \prod_{j\ne i}(e^{\tau_j\lambda}-1)
   =\lambda^r u_i(\lambda),
 \qquad u_i\in A^\times.
\tag{3.7e}\label{eq:coface-unit-row}
\]
It follows exactly that the coface ideal is \((\lambda^r)\), its quotient
has length \(r\), and its Koszul row is chain-equivalent to
\((\lambda^r,0,\ldots,0)\).  Thus contact order and the unframed derived
fat point are independent of the nonzero source frame even though the
units retain its coefficients.

For the generic order-\(r\) multiple elliptic Gamma source, the same
integer is now characterized without choosing one of those models:
\[
\begin{aligned}
 \text{DA transverse depth}
 &=\text{least annihilator degree}
  =\text{arc contact order}\\
 &=\text{fat-point length}
  =\text{period-incidence level}=r.
\end{aligned}
\tag{3.7g}\label{eq:multiple-gamma-order-law}
\]
The equality is a positive source law.  It determines constraint order
exactly; the unit/cocycle calculation shows why order alone is not a
conservative orientation reporter.

Genericity is essential and the failure is itself an augmentation law.
For a genuine resonant period triple
\(\tau_c=\tau_a+\tau_b\), the translation identity
\(T_c=T_aT_b\) gives
\[
 Q=(1+x_a)(1+x_b)-(1+x_c).
\tag{3.7h}\label{eq:resonant-operator}
\]
Its least augmentation degree is one, but the physical arc cancels the
entire linear packet and leaves
\[
              Q\bigl(e^{\tau\lambda}-1\bigr)
                 =\tau_a\tau_b\lambda^2.
\tag{3.7i}\label{eq:resonant-tangency}
\]
Thus contact order is two and tangency excess is one in a natural multiple
elliptic Gamma source.  Source height always bounds contact order from
below; equality fails exactly when the leading resolved monomials cancel
after arc augmentation.  This is a local constraint-theoretic shadow of
\(\ker\Sigma_n\), not an exception to the parent law.

This loss has an exact orientation law.  For
\(G_\tau(\lambda)=\prod_i(e^{\tau_i\lambda}-1)\),
\[
 G_{-\tau}(\lambda)
   =(-1)^{r+1}e^{-(\sum_i\tau_i)\lambda}G_\tau(\lambda).
\tag{3.7j}\label{eq:source-reversal-cocycle}
\]
The multiplier is an involutive cocycle.  On the even-cardinality balanced
locus it is one, so opposite source orientations have identical complete
analytic germs, not merely identical divisors.  Idealization, row
reduction, and an unframed Koszul or Sullivan resolution are therefore
nonconservative observations for source orientation.  They provide exact
constraint depth but cannot replace the source presentation, its reversal
cocycle, and a polarized normalization.  This is the local-algebraic
instance of \eqref{eq:source-conservative}: valuation survives while an
orientation-bearing resolved unit lies in the observation kernel.

This blindness is now exact across analytic realizations.  Additive
difference annihilators and multiplicative translate-product annihilators
pull back from the same squarefree normal-crossing ideal.  Hence source
height, arc contact, fat-point length, period-incidence level, and
polarizing half-density form one realization-independent physical package.
That package is frame-blind: all nonvanishing period frames of one order
give the same values.  The bounded symbol spectrum, by contrast,
reconstructs the embedded period module exactly.  The conservative source
invariant is therefore the pair
\[
        (\text{normal-crossing physical package},
         \text{bounded symbol spectrum}),
\tag{3.7k}\label{eq:physical-package-symbol-pair}
\]
not either marginal.  This is the hypertranscendental realization of the
network law: one quotient retains order while its propagated kernel contains
frame orientation; the second coordinate restores separation before loss.

\paragraph{Prime dilation and the Gamma/co-Poisson cell.}
The half-density prime character
\[
                         b_p(s)=p^{\,s-1/2}
\]
is the common local transport behind the Burnol and Gamma descriptions.
Its removable divided difference is exactly the finite-log
Gamma/co-Poisson interval cell:
\[
 \frac{b_p(s)-b_p(\rho)}{s-\rho}
   =\Gamma_{\log p}(s,\rho),
\tag{3.8}\label{eq:prime-gamma-cell}
\]
with the derivative value at \(s=\rho\).  At any zero \(\rho\) of a
holomorphic function \(Z\), the exact resolvent identity is
\[
 \bigl(b_p(s)-b_p(\rho)\bigr)
       \frac{Z(s)-Z(\rho)}{s-\rho}
 =Z(s)\Gamma_{\log p}(s,\rho).
\tag{3.9}\label{eq:prime-resolvent}
\]
The reflected scale-seam charge is likewise
\[
 \frac{\log p}{\sqrt p}
 \left(b_p(1-\overline\rho)-b_p(\rho)\right).
\tag{3.10}\label{eq:prime-seam}
\]
Thus Gamma/co-Poisson transport, prime dilation, and reflected orientation
are not parallel ingredients: they are divided, transported, and reflected
forms of one half-density character.

\subsection{Compact recurrence and split transport are independent axes}

The polar decomposition \eqref{eq:character-polar} prevents recurrence
rank from being confused with Cartan slope.  For a nonempty finite energy
family with a chosen reference level, let \(r_{\rm temp}\) be the rational
rank of its anchored energy-gap span and let \(r_{\rm rel}\) be the rank of
the exact balanced projective relations.  Rank--nullity gives the
reference-independent conservation law
\[
                         r_{\rm temp}+r_{\rm rel}=N-1.
\tag{3.11}\label{eq:recurrence-rank}
\]
For mutually projectively independent disconnected components, temporal
ranks add.  For a connected finite interaction graph of total tensor
dimension \(D\), the temporal rank is \(D-1\); merging components of
dimensions \(D_A,D_B\) creates exactly
\[
                         (D_A-1)(D_B-1)
\tag{3.12}\label{eq:merge-curvature}
\]
new temporal directions.  These are compact/projective character counts.
Every finite frequency family admits an integer basis, and its phase-orbit
closure is the image of an integer homomorphism between compact tori; for a
rationally independent basis the phase orbit is dense.  Completed observer
time is therefore a compact connected torus whose dimension is rational
frequency rank.  Uniform clock-time sampling converges to its Haar law,
and its identically-one characters are exactly the resonant, zero-detuning
ones.  A nonzero exact return exists exactly when completed-time dimension
is at most one.  Restricting an observer projects completed time onto the
restricted torus and cannot increase temporal dimension.  These theorems
contain no operation on the real exponential factor \(e^{t\mu}\).

The arithmetic source exhibits the separation internally.  A packet of
\(N\) distinct primes has normal rank \(N\), minimal lossless valuation
depth \(N\), and recurrence degree \(N\).  For every spectral point in the
critical strip, independently of those integer identities, the sign of
reflected entropy selects one of the two cofinal Hellinger--Rees species.
Thus arbitrarily rich recurrence or tomography does not force zero Cartan
slope.  It certifies how much compact/source information is present; the
split support question remains separate.

\subsection{Finite common reporting is already source-exact}

One possible loophole is that finite positivity estimates might use a
different witness at each spectral node.  The finite Fisher construction
closes that loophole.  Let \(p_i\ge0\), \(\sum_i p_i=1\), let
\(h_i\ge0\), and let \(|d_i|\le1\).  Put
\[
 \mathcal E_F=
 \sum_{i,j}p_ip_j|h_i-h_j|
       (h_i+h_j-2m)^2.
\tag{3.13}\label{eq:fisher-holonomy-energy}
\]
There is one literal reporter \(R\in L^2(0,\infty)\), independent of the
Hardy node \(w\), such that
\[
                         \|R\|_2^2\le\frac12\mathcal E_F
\tag{3.14}\label{eq:common-reporter-energy}
\]
and, for every \(\Rea w>0\),
\[
 \mathcal L R(w)=
 \sum_{i,j}
 D_{h_j,h_i}\!\left(
 p_ip_j(h_i+h_j-2m)(d_i+d_j);w\right),
\tag{3.15}\label{eq:common-reporter-column}
\]
where
\[
 D_{a,b}(k;w)=
 \begin{cases}
 k\bigl(e^{-wa}-e^{-wb}\bigr)/w,&a\le b,\\
 0,&a>b.
 \end{cases}
\]
The proof assembles the joined codiagonal as a finite sum of interval
indicators.  Its pointwise square is bounded by half the coarea density;
integration gives \eqref{eq:common-reporter-energy}, and Laplace
transformation of every interval gives its semigroup divided difference.

Thus finite common testing is not the missing source.  What remains on
this route is a presentation-independent infinite ordinary-prime
Gamma/co-Poisson reporter mate retaining affine terms and positive
four-unit energy through completion.  No such infinite theorem is claimed
here.

\section{Arithmetic realization of the modular rate}

\subsection{Exact orientation conservation}

For \(s\in\C\), define the reflected shell entropy
\[
 \cE(s)=
 \frac{\log|\chi_2(s)|-\log|\chi_2(1-\overline{s})|}
      {\log2}.
\tag{4.1}\label{eq:entropy}
\]
Directly from \(\chi_2(s)=2^{1/2-s}\),
\[
                         \cE(s)=-\kappa(s).
\tag{4.2}\label{eq:entropy-kappa}
\]

Let \(\Xi\) denote the completed Riemann function.  If
\(m_\rho\) is its analytic order at \(\rho\), set
\[
                    Q_\rho=-m_\rho\left(\rho-\frac12\right).
\tag{4.3}\label{eq:charge}
\]

\begin{theorem}[Riemann realization of Cartan slope]
\label{thm:main}
The following statements hold unconditionally.

\begin{enumerate}[label=\textup{(\roman*)},leftmargin=2em]
\item The rational two-endpoint depth nearby cycle and the completed Xi
      affine logarithmic current have the same outer charge:
      \[
      \begin{aligned}
       w^{-1}\!\left(N_\rho(w^{-1})-1\right)&\longrightarrow Q_\rho,
                 &&w\to0,\ w\ne0,\\
       (s-\rho)\,\cG_\Xi(s)&\longrightarrow Q_\rho,
                 &&s\to\rho,\ s\ne\rho.
      \end{aligned}
      \tag{4.4}\label{eq:outer-inner}
      \]
\item Its horizontal component is the reflected source entropy:
      \[
       \Rea Q_\rho
       =\frac{m_\rho}{2\log2}
          \left(\log|\chi_2(\rho)|
          -\log|\chi_2(1-\overline{\rho})|\right).
      \tag{4.5}\label{eq:charge-entropy}
      \]
      Hence orientation is transported without squaring.
\item The additive duality--adjunction fracture is
      \[
                    \cD(\rho)=\kappa(\rho),
      \qquad
                    \cD(\rho)=0
                    \Longleftrightarrow
                    \Rea\rho=\frac12.
      \tag{4.6}\label{eq:fracture}
      \]
\end{enumerate}
\end{theorem}

\begin{proof}
The two limits in \eqref{eq:outer-inner} both evaluate to
\(-m_\rho(\rho-\tfrac12)\): the first by expansion of the rational depth
successor in the inverse coordinate, the second by the residue of
\(-(\,s-\tfrac12\,)\Xi'(s)/\Xi(s)\).
Equation \eqref{eq:charge-entropy} follows from
\[
 \log|\chi_2(s)|-\log|\chi_2(1-\overline{s})|
 =-2\left(\Rea s-\frac12\right)\log2.
\]
The two nearby-pole branches differ by
\((\rho-\tfrac12)-(\tfrac12-\overline{\rho})=\kappa(\rho)\),
which proves \eqref{eq:fracture}.
\end{proof}

The theorem identifies the arithmetic value of the abstract slope.
Theorem~\ref{thm:cartan-fate} then classifies its possible downstream
fates.  Off-axis reflection is exactly conservative in a hyperbolic plane,
but it cannot carry a nonzero positive conservative metric or even a finite
invariant upper gauge for one surviving nonzero eigenfunctional.

\subsection{Green factorization}

The finite divisor Green operator factors as \(G=-JC_R\), where \(C_R\)
is conjugate Riesz representation and \(J\) is the odd Hodge operator.
Its two transported covariance defects are
\[
 D_J=(\Theta J-J\Theta)C_R,\qquad
 D_R=J(\Theta C_R-C_R\Theta_c).
\tag{4.7}\label{eq:green-defects}
\]
At a divisor copy they are opposite multiples of \(\kappa(\rho)\):
\[
                         (D_J,D_R)=(u,-u).
\tag{4.8}\label{eq:green-antidiagonal}
\]
Therefore the composite Green operator is Frobenius-covariant
unconditionally, while
\[
 (D_J,D_R)=0
 \quad\Longleftrightarrow\quad
 \Rea\rho=\frac12.
\tag{4.9}\label{eq:green-rh}
\]
The source entropy, nearby-cycle fracture, and resolved Green defect are
thus the same oriented normal coordinate in three charts.

Let
\[
                     \Delta_G(a)=(D_J(a),D_R(a)).
\tag{4.10}\label{eq:resolved-green-map}
\]
The distinction between resolved and augmented covariance is exact:
\[
 \Sigma_2\Delta_G=0\quad\hbox{unconditionally},
 \qquad
 \bigl(\Delta_G=0\bigr)\Longleftrightarrow\mathrm{RH}.
\tag{4.11}\label{eq:resolved-green-rh}
\]
Indeed, the forward implication evaluates the first coordinate on a
single divisor basis vector and uses \eqref{eq:green-rh}; the reverse
implication makes the Hodge defect vanish at every divisor copy, after
which \eqref{eq:green-antidiagonal} kills the Riesz coordinate as well.

Consequently any source-selected sector
\(S\subset\cH^1\oplus\cH^1\) proves RH as soon as
\[
       \im\Delta_G\subset S,
       \qquad S\cap\ker\Sigma_2=\{0\}.
\tag{4.12}\label{eq:resolved-source-rh}
\]
This is not presented as a decorative equivalence: the rest of the paper
tests concrete source sectors against its two independent requirements.
The finite balanced Hellinger--Gamma sector passes transversality but not
the support inclusion for a nonzero antidiagonal.  Its naive depth
completion loses closed conservativity.  The missing bridge must satisfy
both parts of \eqref{eq:resolved-source-rh} simultaneously.

\subsection{Asymptotic superselection}

At an actual right-half divisor copy, the reflected prime-valuation
sectors are not merely different finite grades.  One is finite or
half-normal; the other is flat beyond every finite Wold--Rees order.
No finite-Rees-moderate operator sews them.  If an operator sews them
exactly, its norm grows beyond every inverse Rees power.  This is the
asymptotic form of Theorem~\ref{thm:cartan-fate}: positive regular sewing of a
nonzero hyperbolic rate is impossible.

\section{The source-native Meyer quotient}

\subsection{Actual Xi spectral inclusion}

Let \(\cS_{\mathrm{M}}\) be the all-weight logarithmic Meyer source space
and let \(\cB_{\mathrm{P}}\) be the literal Poisson boundary range.  Define
\[
                         \cH_{\mathrm M}
                         =\cS_{\mathrm M}/\cB_{\mathrm P}.
\tag{5.1}\label{eq:meyer-cohomology}
\]
This is an actual source quotient; it is not indexed by zeros.

\begin{theorem}[Finite synthesis and infinite admission]
\label{thm:meyer-wall}
The quotient \(\cH_{\mathrm M}\) has the following properties.

\begin{enumerate}[label=\textup{(\roman*)},leftmargin=2em]
\item Every \(\rho\in\Div(\Xi)\) occurs in the dual dilation flow with
      its full analytic multiplicity.  The jet functionals
      \(E_{\rho,0},\ldots,E_{\rho,m_\rho-1}\) are linearly independent.
\item For every finite set of divisor points, the joint map to all jets
      below their multiplicities is surjective.
\item The evaluation map to the unrestricted product of all Xi jets is
      not surjective.  In particular, the constant unit packet is not
      represented by a Meyer class.
\item There is an actual sequence \(\rho_n\in\Div(\Xi)\) with
      \(|\Im\rho_n|\to\infty\) such that every fixed source Mellin
      transform tends to zero on \(\rho_n\).
\item A source can be chosen with value one at each individual
      \(\rho_n\), but for every differential order \(r>0\) its weighted
      Sobolev cost at stage \(n\) is greater than \((n+1)^r\).
\end{enumerate}
\end{theorem}

\begin{proof}
Poisson--Mellin factorization shows that if zeta vanishes through order
\(k\) at \(\rho\), then the \(k\)-th Mellin jet annihilates every literal
Poisson boundary.  These jets therefore descend to
\(\cH_{\mathrm M}\).  Smooth compactly supported logarithmic profiles
separate every finite jet packet.  The associated confluent Wronskian is
the product of nonzero factorials and pairwise divisor differences, which
proves independence and finite surjectivity.

For \(0\le\sigma\le1\), integration by parts gives
\[
 |\tau|^r\,|M_g(\sigma+i\tau)|
 \le
 \int_{\R}\left|
   \frac{d^r}{dt^r}\bigl(e^{\sigma t}g(t)\bigr)
 \right|dt.
\tag{5.2}\label{eq:vertical-estimate}
\]
The Xi divisor is locally finite and unbounded in imaginary height, so
one may choose \(\rho_n\) escaping vertically.  Equation
\eqref{eq:vertical-estimate} makes every fixed source value tend to zero,
excluding the constant unit packet.  Finite surjectivity supplies an
individual unit source at each stage; applying
\eqref{eq:vertical-estimate} to that source yields the stated
\((n+1)^r\) lower bound.
\end{proof}

The sharp frontier inside the quotient is the source-native spectral
synthesis implication
\[
\left[
 E_{\rho,k}(u)=0
 \text{ for every }\rho\text{ and }k<m_\rho
\right]
\Longrightarrow
u=0\text{ in }\cH_{\mathrm M}.
\tag{5.3}\label{eq:synthesis}
\]
It is equivalent to injectivity of the canonical comparison from
\(\cH_{\mathrm M}\) to compact-open Xi cohomology.  It is not required
for the minimal RH implication below, because one-way inclusion and
nonvanishing of every actual Xi evaluation are already proved.

\subsection{Archimedean completion and the common-source resolvent}

There is also a literal Hilbert-space realization before quotient.  Let
\[
 \lambda_\rho=-i\left(\rho-\frac12\right),\qquad
 F(x)=\Xi\!\left(\frac12+ix\right),
\]
let \(A\) be multiplication by the real coordinate on its maximal
\(L^2(\R)\) domain, and let \(q_\rho\) be the removable first divided Xi
function.  Then \(A\) is symmetric, \(F\ne0\), and for every actual Xi
divisor point,
\[
 q_\rho\ne0,\qquad q_\rho\in\dom A,\qquad
                       Aq_\rho=\lambda_\rho q_\rho+F.
\tag{5.4}\label{eq:common-source-resolvent}
\]
The same nonzero vector \(F\) is the defect for every divisor point.
Moreover its full real phase orbit has dense span in \(L^2(\R)\), as does
the translation orbit of the physical theta source.

Equation \eqref{eq:common-source-resolvent} is not a Hilbert--P\'olya
eigenvalue equation.  It is an affine common-source resolvent equation.
Killing \(F\) formally creates all desired eigenvectors, but any closed
phase-invariant radical containing \(F\) contains its dense orbit.  This is
the exact Hedenmalm dense-radical shadow of the singular-descent law: the
source defect is universal, while regular quotienting of that defect
collapses the carrier.

\subsection{Reciprocal flow and weight-one pairing}

The reciprocal operator
\[
                           (Jg)(x)=x^{-1}g(x^{-1})
\tag{5.5}\label{eq:paugam-j}
\]
preserves the all-weight source, carries literal Poisson boundaries to
boundaries, and descends to an involution on \(\cH_{\mathrm M}\).
For a zero jet datum of order \(n\), the flow-correct alternating form is
not the top-jet determinant.  It is the complementary-jet block
\[
\Psi_{\rho,n}(u,v)
=\sum_{j=0}^{n}\binom nj
\left(
 E_{\rho,j}(u)E_{\rho,n-j}(Jv)
-E_{\rho,j}(v)E_{\rho,n-j}(Ju)
\right).
\tag{5.6}\label{eq:paugam-pairing}
\]
It is alternating, anti-invariant under simultaneous \(J\), and obeys
the exact weight-one dilation law
\[
 \Psi_{\rho,n}(U_cu,U_cv)=e^c\Psi_{\rho,n}(u,v).
\tag{5.7}\label{eq:weight-one}
\]
The binomial completion is forced by the triangular translation of
generalized jets.

\subsection{Why the available positivity is in the wrong degree}

At each positive radius, the local Calder\'on defect quotient is one
complex dimension.  A single complex line cannot carry the required odd
real structure, but reciprocal doubling can.  On the doubled defect
quotient there is an anti-linear operator squaring to \(-1\), and its
cup--Hodge diagonal is exactly
\[
                          |z_+|^2+|z_-|^2.
\tag{5.8}\label{eq:local-positive}
\]
This form is positive and has zero radical.

However, every genuine global degree-one Meyer cycle satisfies the
Poisson graph at every radius.  Its local Calder\'on defect class is
therefore zero.  Consequently \eqref{eq:local-positive} pulls back
identically to zero on genuine middle cycles.  Local positivity belongs
to the transverse regulator, or degree-two defect, not to the
degree-one cohomology which carries Xi eigenfunctionals.

This grading obstruction removes the main ambiguity in the frontier:
another positive local block cannot prove RH, regardless of how sharp
its scalar estimate is.

\subsection{A source-explicit co-Poisson--zeta bridge}

The Burnol channel is not merely a downstream positivity consumer.  For
\(0<b<1\), take the reciprocal of the indicator seed of \([b,b^{-1}]\).
Its right co-Poisson sum is, almost everywhere on the positive ray,
\[
 P_b(t)=\frac{\{bt\}-\{t/b\}}{t}.
\tag{5.8a}\label{eq:burnol-fractional-source}
\]
The full critical-strip Mellin transform is exact:
\[
 \int_0^\infty P_b(t)t^{-w}\,dt
   =\bigl(b^{-w}-b^w\bigr)\frac{\zeta(w)}{w},
 \qquad 0<\Rea w<1.
\tag{5.8b}\label{eq:burnol-zeta-mellin}
\]
Because \(0<b<1\), the window factor does not vanish in this strip.
Therefore this one elementary co-Poisson source vanishes in Mellin space
exactly at the zeta divisor.  The Burnol route already has a literal
source-to-divisor transform before quotient positivity is asked; its open
edge is the descended polarized comparison, not spectral relevance.

\subsection{Faithful co-Poisson positivity and the same hyperbolic rate}

On every closed Burnol co-Poisson quotient in the right regularity strip,
the completed Mellin evaluation kernels have dense linear span.  Hence,
for a positive operator supported on that quotient, its Berezin diagonal
\[
                         \cB_T(w)=\langle k_w,Tk_w\rangle
\tag{5.9}\label{eq:berezin}
\]
is faithful: if \(\cB_T(w)=0\) throughout the strip, then \(T=0\).
Cosine reflection transports the left evaluation kernel to the right
kernel at \(1-w\).  Therefore a Fourier-odd operator has reflection-odd
Berezin diagonal.

For the finite bilateral aperture holonomy, the diagonal at a
co-Poisson annihilator point has the exact normalized value
\[
 \frac{\Rea\cB_T(w)}{\|k_w\|^2}
   =r^{\,2\Rea w-1}-r^{\,1-2\Rea w}
   =2\sinh\!\left(\log r\,\kappa(w)\right).
\tag{5.10}\label{eq:burnol-hyperbolic}
\]
Thus the Sonin--Burnol positive consumer does not introduce a new
orientation law.  It reads the same transverse character through its
hyperbolic sine.  Faithfulness says positivity can detect a surviving
operator; reflection oddness says the oriented reader remains signed.

Compression to the actual co-Poisson source quotient preserves
self-adjointness and Fourier oddness.  Consequently its holonomy \(T\)
satisfies the exact operator-order alternative
\[
                         T\ge0\quad\Longleftrightarrow\quad
                         T=0\quad\Longleftrightarrow\quad -T\ge0.
\tag{5.11}\label{eq:burnol-order-flatness}
\]
Moreover its normalized even--odd parity block is twice the cross-Gram of
the actual forward quotient channels; flat holonomy is equivalent to
vanishing of that cross-Gram.  This is the Hilbert-operator shadow of
Theorem~\ref{thm:cartan-fate}: Fourier-odd reciprocal transport survives in
an off-diagonal pairing, but no nonzero one-sided positive order.

Writing \(F_{a,b,r}\) and \(I_{a,b,r}\) for the actual forward and inverse
full quotient channels, the order assertion has the exact factorization
form
\[
 -T\ge0
 \quad\Longleftrightarrow\quad
 \exists C,\qquad \|C\|\le1,\qquad F_{a,b,r}=I_{a,b,r}C.
\tag{5.11a}\label{eq:burnol-douglas}
\]
Fourier oddness then forces every such contractive factorization to have
\(T=0\) and zero parity cross-Gram.  Hence Burnol positivity is not an
unidentified sign estimate: it is precisely the polarized factorization
gate of \eqref{eq:douglas-descent}, and its surviving obligation is exact
channel comparison.

\subsection{Weyl-even conditioning and Weyl-odd holonomy}

The absolute Cartan coordinate is independently forced by two source
geometries.  Write
\[
                         w=\frac12+i\alpha,
             \qquad \mu=\Rea w-\frac12=-\Im\alpha.
\]
For the Hedenmalm Borel character on the window
\(T_j=r^{2j}-1\), the product of the two reflected essential suprema is
exactly
\[
                         C_j=r^{j|\mu|}.
\tag{5.12}\label{eq:borel-condition}
\]
For the Gamma-completed Sonin evaluation vector \(k_{w,a}\), provided the
base vector is nonzero,
\[
 \lim_{N\to\infty}
 \frac{\log\|k_{w,r^{-N}a}\|}{N\log r}=|\mu|.
\tag{5.13}\label{eq:sonin-aperture}
\]
The Borel condition exponent and the Sonin aperture exponent are therefore
one number, in constructions with different source measures and different
completions.  Both vanish exactly on the critical line.

Equation~\eqref{eq:burnol-hyperbolic} is the oriented companion:
\[
 H_r(\mu)=2\sinh(2\log r\,\mu),\qquad
 |H_r(\mu)|=e^{2\log r|\mu|}-e^{-2\log r|\mu|}.
\tag{5.14}\label{eq:odd-even-cartan}
\]
Thus aperture and condition numbers are Weyl-even readers of the Cartan
slope, while reflected holonomy is a Weyl-odd reader.  They cannot disagree
about criticality because they are functorial transforms of the same
character, not independent positivity mechanisms.

\clearpage
\section{Shadows of polarized exact descent}

The following rows are not separate main theorems.  Each discharges or
obstructs separation, gluing, stability, or Hermitian calibration in
Theorem~\ref{thm:exact-descent}.

\begin{center}
\small
\begin{tabular}{@{}>{\raggedright\arraybackslash}p{0.20\linewidth}
                >{\raggedright\arraybackslash}p{0.31\linewidth}
                >{\raggedright\arraybackslash}p{0.39\linewidth}@{}}
\toprule
Observed fate & Source-sector mechanism & Arithmetic realization\\
\midrule
Oriented conservation
 & equivariant regular transport preserves signed slope
 & shell entropy, Xi charge, nearby cycle, depth fracture, and resolved Green defects equal \(\kappa\)\\\addlinespace[2pt]
Positive descent
 & invariant seminorm retains a vector only at slope zero; otherwise it enters the radical
 & half-density continuous spectrum and Wold quadrant boundedness select the critical line\\\addlinespace[2pt]
Closed graph
 & a nonzero dual slope produces capacity \(e^{-n\mu^2}\) with fixed output, hence full scalar verticality
 & Meyer Xi graph closure is single-valued on line and maximally vertical off line\\\addlinespace[2pt]
Two-sided observation
 & reflected condition growth has exponent \(|\mu|\)
 & Borel condition numbers and Sonin aperture norms have the same exact exponent\\\addlinespace[2pt]
Reciprocal pairing
 & opposite slopes retain only a cross form; positive invariant form vanishes
 & Green antidiagonal, Paugam weight-one pairing, and Fourier-odd Burnol holonomy\\\addlinespace[2pt]
Singular retention
 & resolved or strictly graded data retain what regular descent erases
 & finite-fibre Rees carrier exists exactly at local finite depth; off-line sectors are flat beyond every finite order\\\addlinespace[2pt]
Exact gluing
 & compatible observations have a unique source lift
 & normal-crossing face descent, relative Sullivan attachment, and the safe-chamber M\"obius--Gaussian Neumann source\\\addlinespace[2pt]
Descent order
 & total tests uniquely fix their reporter, so strict energy reduction requires changed carrier or incidence before polarization
 & the full Euler damping packet forbids every same-carrier compression from eight units to four\\\addlinespace[2pt]
Hermitian calibration
 & Gram order is equivalent to contractive factorization; topology fixes only a conormal line
 & full Euler Fisher energy is four, division-free reporter energy is eight, and Burnol quotient positivity is Douglas factorization\\\addlinespace[2pt]
Compact recurrence
 & projective phase rank plus relation rank is dimension minus global phase
 & recurrence rank, graph curvature, and prime-normal rank do not determine entropy sign\\\addlinespace[2pt]
Higher-rank comparison
 & bounded prefix error preserves normalized Cartan vectors
 & finite differential embeddings with pseudodifferential right inverse preserve HT Cartan spectra\\
\bottomrule
\end{tabular}
\end{center}

Additional source theorems determine which row is available.  Shell and
theta systems prove local divisor fullness; Euler--Hellinger analysis
separates four-unit capacity from the eight-unit division-free reporter;
the finite common Hardy reporter proves that finite source-to-test
realization is literal rather than nodewise; Hedenmalm cyclicity makes
radicalization dense; local Calder\'on positivity lives in regulator degree
and vanishes on genuine middle cycles.  None is a second explanation.  Each
constructs a source sector, supplies compatible lifts, tests separation or
stability, or supplies a polarized consumer.

\section{The exact RH frontier: ordered polarized singular arithmetic descent}

\subsection{Minimal pointed criterion}

Let \(U_c\) be dilation on \(\cH_{\mathrm M}\) and define the centered
flow
\[
                         V_c=e^{-c/2}U_c.
\tag{7.1}\label{eq:centered-flow}
\]
At a Xi divisor point, the zero-order functional transforms by the
character \(e^{c(\rho-1/2)}\).

\begin{theorem}[Target-free Gaussian frontier]
\label{thm:fixed-source}
The logarithmic Gaussian
\[
                 g_\ast(t)=e^{-\pi t^2}
\]
satisfies the full all-weight Meyer source condition and has entire Mellin
transform
\[
                 M_{g_\ast}(s)=e^{s^2/(4\pi)}\ne0.
\]
Hence its class \(u_\ast=[g_\ast]\in\cH_{\mathrm M}\) satisfies
\[
                         E_{\rho,0}(u_\ast)\ne0
\]
for every actual Xi divisor point.  For this one class,
\[
\left[
 \forall\rho\in\Div(\Xi)\;\exists C_\rho\in\R\;
 \forall c\in\R,\quad |E_{\rho,0}(V_cu_\ast)|\le C_\rho
\right]
\quad\Longleftrightarrow\quad\mathrm{RH}.
\tag{7.2}\label{eq:fixed-source-rh}
\]
\end{theorem}

\begin{proof}
The Gaussian transform is an exponential and therefore has no zero; the
Meyer evaluation identity gives nonvanishing under every actual zeroth
evaluation.  Its centered orbit satisfies
\[
 |E_{\rho,0}(V_cu_\ast)|
 =e^{c(\Rea\rho-1/2)}|E_{\rho,0}(u_\ast)|.
\]
If the orbit is bounded for both signs of \(c\), its nonzero initial value
forces \(\Rea\rho=\tfrac12\).  Conversely, on the critical line the orbit
has constant modulus.  Apply this pointwise to the entire Xi divisor.
\end{proof}

The fixed-source formulation has a reflection-adapted strengthening.  Put
\[
 h(t)=e^{-t/2}e^{-\pi t^2},\qquad H_c=V_ch+V_{-c}h,
\tag{7.3}\label{eq:reflected-countertransport}
\]
and write the same symbols for their Meyer classes.

\begin{theorem}[Uniform-capacity reflected separation]
\label{thm:reflected-capacity}
For every \(c\in\R\), \(JH_c=H_c\) and
\[
                         \overline p_{1/2}(H_c)\le2.
\tag{7.4}\label{eq:uniform-capacity}
\]
For every actual Xi divisor point \(\rho\),
\[
 E_{\rho,0}(H_c)=
 \left(e^{c(\rho-1/2)}+e^{-c(\rho-1/2)}\right)
 \exp\!\left(\frac{(\rho-1/2)^2}{4\pi}\right).
\tag{7.5}\label{eq:reflected-evaluation}
\]
Consequently,
\[
 \sup_{c\in\R}|E_{\rho,0}(H_c)|<\infty
 \quad\Longleftrightarrow\quad
 \Rea\rho=\frac12,
\tag{7.6}\label{eq:reflected-selector}
\]
and uniform boundedness at every actual Xi point is equivalent to RH.
\end{theorem}

\begin{proof}
Weighted reflection reverses centered flow, while \(Jh=h\); hence
\(JH_c=H_c\).  Direct Gaussian integration gives
\(p_{1/2}(h)=1\).  Centered flow is isometric, so the triangle inequality
gives \eqref{eq:uniform-capacity}, also after quotienting.  The centered
Mellin character and the Gaussian Mellin transform give
\eqref{eq:reflected-evaluation}.  Its Gaussian coefficient never vanishes.
If \(\Rea\rho=\tfrac12\), both exponentials have modulus one.  If
\(\Rea\rho\ne\tfrac12\), one term grows exponentially in one time
direction and cannot be cancelled by its reciprocal term.  This proves
\eqref{eq:reflected-selector}; applying it to every actual divisor point
gives the last assertion.
\end{proof}

There is an explicit Green sequence behind the same boundary.  For
\(a\ne0\), define
\[
 q_a=e^{-a^2/(4\pi)},\qquad
 G_{a,N}(t)=\sum_{r=1}^{N}r^{-1/2}q_a^r e^{-\pi t^2/r}.
\tag{7.6a}\label{eq:gaussian-green-partial}
\]
Every term is positive and direct Gaussian integration gives
\[
 \int_\R G_{a,N}(t)\,dt=\sum_{r=1}^{N}q_a^r
 \le \frac{q_a}{1-q_a},
 \qquad
 \int_\R e^{at}G_{a,N}(t)\,dt=N.
\tag{7.6b}\label{eq:gaussian-green-defect}
\]
Undoing half-density and adding any fixed vertical phase produces genuine
Meyer source partial sums with uniformly bounded half-density mass, while
their Mellin values at horizontal displacement \(a\) diverge linearly.
Thus off-line failure is not lack of source mass or a weak estimate: one
explicit positive Green primitive remains bounded in the completed source
norm and accumulates an unbounded singular evaluation.  In
Theorem~\ref{thm:augmentation-conservativity} this is quantitative loss of
completed transversality; regular completion cannot turn the corresponding
source sector into a uniformly conservative one.

The network law also determines the exact singular compensation, not only its
existence.  Fix \(c\ne\rho\) and define
\[
 D_{c,\rho}(s)=
   \exp\!\left(\frac{(\rho-c)^2}{4\pi}\right)
   -\exp\!\left(\frac{(s-c)^2}{4\pi}\right),
 \qquad
 F_{c,\rho}(s)=\frac1{\zeta(s)D_{c,\rho}(s)}.
\tag{7.6c}\label{eq:singular-gaussian-lift}
\]

\begin{proposition}[Forced local singularity budget]
\label{prop:singular-budget}
If \(\rho\) is a Xi zero of multiplicity \(m_\rho\) in the open critical
strip, then \(D_{c,\rho}\) has a simple zero at \(\rho\), while
\[
             \operatorname{ord}_\rho F_{c,\rho}=-(m_\rho+1).
\tag{7.6d}\label{eq:successor-pole}
\]
On a punctured neighbourhood of \(\rho\),
\[
 \zeta(s)F_{c,\rho}(s)=D_{c,\rho}(s)^{-1},
 \qquad
 \zeta(s)D_{c,\rho}(s)F_{c,\rho}(s)=1.
\tag{7.6e}\label{eq:singular-lift-anomaly}
\]
Thus a global Poisson--Gaussian source lift cannot be regular at an actual
divisor point: it must assemble the successor principal part of order
\(m_\rho+1\).
\end{proposition}

\begin{proof}
The derivative of \(D_{c,\rho}\) at \(\rho\) is
\[
 -\exp\!\left(\frac{(\rho-c)^2}{4\pi}\right)
       \frac{\rho-c}{2\pi}\ne0,
\]
so its order is one.  In the open strip, \(\zeta\) and Xi have the same
order \(m_\rho\).  Orders add under multiplication and change sign under
inversion, proving \eqref{eq:successor-pole}.  The punctured identities are
ordinary cancellation where both factors are nonzero.
\end{proof}

The singular lift already has a literal arithmetic source on its safe
chamber.  Let \(\chi_c\) be the Gaussian Mellin character and put
\[
 r_{c,\rho}(s)=\chi_c(s)\chi_c(\rho)^{-1}.
\]
If \(\Rea s>1\) and \(|r_{c,\rho}(s)|<1\), then the M\"obius--Gaussian
Neumann series is absolutely summable and
\[
 \sum_{k\ge0}
   L\!\left(m_1(\mu),s\right)\chi_c(\rho)^{-1}
       r_{c,\rho}(s)^k
   =F_{c,\rho}(s).
\tag{7.6f}\label{eq:mobius-gaussian-neumann}
\]
Here \(L(m_1(\mu),s)=\zeta(s)^{-1}\) on the same half-plane.  Thus neither
the arithmetic coefficients nor the Gaussian convolution-depth source is
missing.  The open problem begins exactly when this coefficient-faithful
section must be continued from the regular chamber to Xi-zero fibres in the
nonlocal singular topology while retaining its polarized source lift.

Even the local normalization at the singular boundary is forced.  Writing
\(A_\infty\) for the archimedean Xi factor and \(\chi_c\) for the Gaussian
Mellin character, the punctured completed residue is
\[
 \lim_{s\to\rho}(s-\rho)\,\Xi(s)F_{c,\rho}(s)
   =-\frac{2\pi A_\infty(\rho)}
           {\chi_c(\rho)(\rho-c)}\ne0.
\tag{7.6g}\label{eq:gamma-green-residue}
\]
It is independent of the multiplicity of \(\rho\).  Hence a source-faithful
continuation cannot multiply local principal parts by arbitrary scalars:
the Gamma/Green boundary coefficient is already fixed and nonzero.  What
remains is simultaneous global assembly on the descended carrier, not a
choice of local residue.

\begin{definition}[Pointwise invariant evaluation gauge]
\label{def:gauge}
For \(\rho\in\Div(\Xi)\), an admissible gauge is a finite function
\[
                         g_\rho:\cH_{\mathrm M}\longrightarrow[0,\infty)
\]
such that, for some \(C_\rho\ge0\),
\[
 g_\rho(V_cu)=g_\rho(u),\qquad
 |E_{\rho,0}(u)|\le C_\rho g_\rho(u)
\tag{7.7}\label{eq:gauge}
\]
for every \(c\in\R\) and \(u\in\cH_{\mathrm M}\).
\end{definition}

\begin{theorem}[Exact pointwise tempered-admission frontier]
\label{cor:rh}
For an actual Xi divisor point \(\rho\), an admissible gauge exists if
and only if \(\Rea\rho=\tfrac12\).  Consequently,
\[
 \bigl(\forall\rho\in\Div(\Xi)\;\exists g_\rho
       \text{ satisfying \eqref{eq:gauge}}\bigr)
 \quad\Longleftrightarrow\quad \mathrm{RH}.
\tag{7.8}\label{eq:gauge-rh}
\]
\end{theorem}

\begin{proof}
The evaluation is nonzero and satisfies
\[
              E_{\rho,0}V_c
              =e^{c(\rho-1/2)}E_{\rho,0}.
\]
Theorem~\ref{thm:cartan-fate}, with \(z=\rho-\tfrac12\), proves the forward
implication.  If \(\Rea\rho=\tfrac12\), take
\(g_\rho(u)=|E_{\rho,0}(u)|\); the same character law makes this gauge
invariant, with \(C_\rho=1\).  The global equivalence follows because the
Meyer quotient contains every actual Xi zero.
\end{proof}

\subsection{Canonical half-density tempered geometry}

For \(g\in\cS_{\mathrm M}\), define
\[
 p_{1/2}(g)=\int_{\R}e^{t/2}|g(t)|\,dt
\tag{7.9}\label{eq:half-density}
\]
and descend it through the literal Poisson boundary by
\[
 \overline p_{1/2}([g])
   =\inf_{b\in\cB_{\mathrm P}}p_{1/2}(g+b).
\tag{7.10}\label{eq:quotient-half-density}
\]

\begin{theorem}[Exact half-density continuity selector]
\label{thm:half-density}
The following statements hold unconditionally.
\begin{enumerate}[label=\textup{(\roman*)},leftmargin=2em]
\item \(p_{1/2}\) is a seminorm on the strong logarithmic Meyer source and
      \(\overline p_{1/2}\) is its quotient seminorm on
      \(\cH_{\mathrm M}\).
\item Centered dilation is exactly isometric:
      \[
                  \overline p_{1/2}(V_cu)=\overline p_{1/2}(u).
      \]
\item The Hausdorff quotient by the seminorm kernel has a Banach
      completion on which \(V_c\) extends to a surjective linear isometry
      with inverse \(V_{-c}\).
\item For every actual Xi divisor point,
      \[
      E_{\rho,0}\text{ is bounded by }\overline p_{1/2}
      \quad\Longleftrightarrow\quad
      \Rea\rho=\frac12.
      \tag{7.11}\label{eq:selector}
      \]
      On the critical line the bound is contractive.
\item Consequently,
      \[
      \mathrm{RH}
      \quad\Longleftrightarrow\quad
      \text{every actual \(E_{\rho,0}\) is
      \(\overline p_{1/2}\)-bounded}.
      \tag{7.12}\label{eq:canonical-rh}
      \]
\end{enumerate}
\end{theorem}

\begin{proof}
The change of variables \(u=t-c\) gives
\[
\begin{aligned}
p_{1/2}(V_cg)
 &=\int_{\R}e^{t/2}e^{-c/2}|g(t-c)|\,dt\\
 &=\int_{\R}e^{u/2}|g(u)|\,du=p_{1/2}(g).
\end{aligned}
\]
Taking the quotient infimum and using both \(c\) and \(-c\) gives exact
quotient invariance.  Quotient by the null space and completion therefore
extend the centered flow isometrically; the inverse flow proves
surjectivity.

If \(\Rea\rho=\tfrac12\), the Mellin estimate gives
\[
 |E_{\rho,0}([g])|
 \le p_{1/2}(g+b)
\]
for every boundary representative \(g+b\).  Taking the infimum proves
the contractive bound.  Conversely, \(\overline p_{1/2}\) is an admissible
gauge in Definition~\ref{def:gauge}, so Theorem~\ref{cor:rh} forces
\(\Rea\rho=\tfrac12\).  This proves \eqref{eq:selector}; applying it to
every actual Xi point gives \eqref{eq:canonical-rh}.
\end{proof}

\subsection{Exact graph-closure phase transition}

Let \(q:\cH_{\mathrm M}\to\widehat{\cH}_{1/2}\) be the canonical map to
the Hausdorff half-density completion and let
\[
 \cR_\rho=
 \{(q(u),E_{\rho,0}(u)):u\in\cH_{\mathrm M}\}
 \subset\widehat{\cH}_{1/2}\times\C.
\tag{7.13}\label{eq:xi-graph}
\]

\begin{theorem}[Closed Xi graph alternative]
\label{thm:xi-graph-closure}
For every actual Xi divisor point \(\rho\),
\[
 \left[(0,y)\in\overline{\cR_\rho}\Longrightarrow y=0
       \text{ for every }y\in\C\right]
 \quad\Longleftrightarrow\quad
 \Rea\rho=\frac12.
\tag{7.14}\label{eq:graph-single-valued}
\]
If \(\Rea\rho\ne\tfrac12\), then the stronger identity holds:
\[
                         \{0\}\times\C
                         \subset\overline{\cR_\rho}.
\tag{7.15}\label{eq:graph-full-vertical}
\]
Hence every closed linear relation containing the full algebraic Xi graph
is maximally vertical over zero, and no closed single-valued extension of
that evaluation exists on the full completion image.
\end{theorem}

\begin{proof}
On the critical line, Theorem~\ref{thm:half-density} supplies a continuous
extension of \(E_{\rho,0}\); its closed graph has zero vertical fibre and
contains \(\cR_\rho\).

Off line, put \(z=\rho-\tfrac12\), \(\mu=\Rea z\), choose the logarithmic
Gaussian class \(u_\ast\), and set
\[
 u_n=e^{-t_nz}V_{t_n}u_\ast,\qquad t_n=n\mu.
\tag{7.16}\label{eq:xi-graph-null}
\]
Theorem~\ref{thm:cartan-fate} gives
\[
 E_{\rho,0}(u_n)=E_{\rho,0}(u_\ast)\ne0,\qquad
 \|q(u_n)\|=e^{-n\mu^2}\|q(u_\ast)\|\longrightarrow0.
\]
Thus \((0,E_{\rho,0}(u_\ast))\in\overline{\cR_\rho}\).  Scalar closure
and nonvanishing give \eqref{eq:graph-full-vertical}.  The equivalence and
the assertion about every larger closed relation follow immediately.
\end{proof}

\begin{corollary}[Algebraic--tempered spectral split]
\label{cor:spectral-split}
Let
\[
\Spec_{\rm alg}^{\Xi}(\cH_{\mathrm M})=\Div(\Xi),\qquad
\Spec_{\rm temp}^{\Xi}(\cH_{\mathrm M})
=\{\rho\in\Div(\Xi):E_{\rho,0}
  \text{ is }\overline p_{1/2}\text{-bounded}\}.
\]
Then
\[
\Spec_{\rm temp}^{\Xi}(\cH_{\mathrm M})
=\Div(\Xi)\cap\{\Rea s=\tfrac12\},
\qquad
\mathrm{RH}\Longleftrightarrow
\Spec_{\rm alg}^{\Xi}=\Spec_{\rm temp}^{\Xi}.
\]
\end{corollary}

\begin{proof}
This is Theorem~\ref{thm:half-density}(iv)--(v).
\end{proof}

\subsection{What has and has not been reduced}

Theorem~\ref{cor:rh} removes assumptions previously mistaken for the
frontier.  Zero exclusion requires neither a new Hilbert space, a new
positive form, nor surjectivity onto a global divisor model.  Exact logical
datum is resolved source conservativity: every actual Green defect must be
supported in a completed arithmetic sector transverse to the codiagonal
kernel.  Zero-Cartan admission of every Xi functional is its rank-one
continuous-dual form.
Theorem~\ref{thm:half-density} supplies one canonical common seminorm and
its positive isometric completion; Theorem~\ref{thm:xi-graph-closure}
proves that off-line failure is full verticality, not a small domain defect.

Both selectors are exact, so asserting their missing boundedness direction
is equivalent to RH.  Their research content is diagnostic rather than a
disguised proof: any proof must derive boundedness from arithmetic source
structure, not from the unknown locations of zeros.  The corpus has fixed
every surrounding ingredient:

\begin{enumerate}[leftmargin=2.2em]
\item the source space and literal boundary quotient are explicit;
\item every Xi evaluation is already a nonzero algebraic functional;
\item centered dilation, reciprocal reflection, and weight-one alternating
      pairing are explicit;
\item projective normal tomography supplies the spectral coordinate
      before ordinary collapse; and
\item Theorem~\ref{thm:finite-fibre} gives an exact topology constraint;
\item the leading physical Xi jet lands coefficient-faithfully in completed
      Hellinger--Gamma incidence, and its unsquared Hodge trace detects zero
      modular displacement exactly; complete prime-normal residual sources
      retain the same nonzero leading conormal by source-derived prime
      isolation;
\item atomic normal-crossing faces have unique compatible gluing, and a
      relative Sullivan tower gives a surjective quasi-isomorphic resolution
      of the square-zero target;
\item the singular Poisson--Gaussian lift has a literal summable
      M\"obius--Gaussian Neumann source on its absolute-convergence chamber,
      and its completed Gamma/Green residue is fixed and nonzero;
\item one reciprocal-indicator co-Poisson source has Mellin transform
      \((b^{-w}-b^w)\zeta(w)/w\) throughout the critical strip;
\item one reflected-Hodge Gaussian family has unconditional uniform source
      capacity at every transport depth; and
\item one literal finite common Hardy reporter simultaneously realizes all
      finite divided-difference nodes with controlled energy;
\item the direct full-Euler Fisher carrier has four-unit critical energy,
      while the division-free pre-scalar completed reporter has exact limiting
      energy eight and a retained Xi graph factorization; totality of its
      damping orbit forbids every same-carrier four-unit reporter preserving
      the full coefficient packet; and
\item Hilbert Gram order is exactly contractive right factorization, including
      the actual Burnol forward/inverse quotient channels.
\end{enumerate}

The exact-descent audit separates what remains into the three gates of
Theorem~\ref{thm:exact-descent}.  Algebraic local gluing is no longer
missing: face descent, physical/Hodge incidence, and the Neumann source
provide it in their respective chambers.  Topological continuation is not
an ordinary completion: informative prime-normal topology must be strictly
finer than product observation, reflected sewing is beyond every finite
Rees order, and Proposition~\ref{prop:singular-budget} forces a successor
pole of order \(m_\rho+1\).  Hermitian calibration is independent and
ordered: the source has a four-unit carrier, the division-free completed
reporter has energy eight, its total coefficient packet makes same-carrier
compression impossible, and projective topology cannot select a factor-two
normalization.  A successful route must therefore identify the exact
compatibility quotient first and polarize only after descent.

The first unresolved object is therefore one \emph{ordered polarized
singular exact descent}: a global source-derived comparison from
coefficient-faithful physical/Hodge incidence to the actual Green resolved
range which (i) identifies and quotients precisely the compatibility-exact
part of the total pre-scalar packet, (ii) glues every surviving compatible
arithmetic packet uniquely, (iii) assembles the forced principal parts with
their fixed Gamma/Green residues without losing conservativity in the
nonlocal completion, and only then (iv) realizes the required Hermitian Gram
order, equivalently the appropriate contractive factorization, while
retaining unsquared orientation.  At the resolved level its separation and
stability consequences include
\[
 \im\Delta_G\subset S_{\rm ar},
 \qquad
 \overline{S_{\rm ar}}\cap\ker\Sigma_2=\{0\},
\tag{7.17}\label{eq:missing-resolved-sector}
\]
for one source-derived arithmetic sector whose closure is taken in the
actual Gamma/co-Poisson completion.  Its gluing gate is the corresponding
surjectivity onto the compatible incidence sector, and its polarized gate
is a source-derived Gram domination, not an assumed downstream positivity
inequality.  On the Hodge incidence route, the same object may instead be
discharged by a source theorem forcing the trace in
\eqref{eq:theta-hodge-detector} to vanish after its displacement has been
identified with the Green coordinate.  Its minimal scalar form is the
fixed-source orbit estimate in \eqref{eq:fixed-source-rh}; its
reflection-adapted form is uniform evaluation of the capacity-bounded
family in Theorem~\ref{thm:reflected-capacity}.  Its canonical common form is
\[
  \forall\rho\in\Div(\Xi),\quad
  |E_{\rho,0}(u)|\le C_\rho\overline p_{1/2}(u)
  \quad(u\in\cH_{\mathrm M}),
\tag{7.18}\label{eq:missing-admission}
\]
with finite \(C_\rho\), derived from ordinary-prime/Gamma/co-Poisson source
incidence without assuming \(\Rea\rho=\tfrac12\).  Each form is equivalent
to RH; the last additionally places the complete actual Xi divisor in one
canonical continuous dual.

The following substitutes are now excluded.

\begin{enumerate}[leftmargin=2.2em]
\item A regular comparison from the complete \(\ell^1\) source erases
      reflection displacement.  It can even remain pointwise nonzero at
      every Xi divisor while identifying all source labels.
\item A closed ungraded relation retaining the dense prime-normal family
      becomes \(\dom\cR\times Y\).
\item More sharply, the closed half-density graph of each off-line actual
      Xi evaluation contains \(\{0\}\times\C\); taking closure cannot repair
      admission.
\item An unrestricted product of divisor jets is not Meyer-admissible.
\item Uniform finite-order cofinal synthesis is impossible.
\item Finite-Rees sewing cannot join an off-line reflected pair.
\item A same-carrier four-unit Euler reporter cannot preserve the complete
      division-free damping packet; totality forces the eight-unit reporter.
\item Local reciprocal Calder\'on positivity vanishes on every genuine
      middle cycle.
\item Green covariance and scalar positivity after codiagonal cannot see
      the antidiagonal defect.
\end{enumerate}

Global spectral synthesis \eqref{eq:synthesis} is a natural stronger
theorem: it would identify all of compact-open Xi cohomology with the
Meyer quotient.  It is not required for
Theorem~\ref{thm:fixed-source} or Theorem~\ref{cor:rh}.  The literal minimal
frontier is one Gaussian orbit bound; \eqref{eq:missing-admission} is the
rank-one source-geometric form of the resolved-sector assertion
\eqref{eq:missing-resolved-sector}.

\subsection{Why several routes converge here}

Shell, theta, Euler--Gamma, Burnol--Sonin, Meyer, recurrence, and
differential-operator constructions are genuinely different sources and
transports.  They are unified by invariant, not by forcing them into one
artificial chain.

\begin{itemize}[leftmargin=2em]
\item The shell route gives the canonical normal coordinate and exact
      finite-fibre topology boundary.
\item The theta route proves that all local Xi jets arise from one
      target-free physical source.
\item The Euler route separates four-unit Fisher capacity from an
      eight-unit division-free reporter; its finite Hellinger--Gamma sector
      is conservative, while naive depth completion verticalizes.
\item Normal-crossing face descent and the relative Sullivan tower supply
      exact algebraic gluing, while showing why target exactness alone does
      not retain source presentation.
\item The Calder\'on--Witt route supplies a nonzero ordered arithmetic
      cross-effect before ordinary cup augmentation kills its antidiagonal.
\item The Burnol--Sonin/Fisher route supplies exact even and odd Cartan
      readers, one finite common reporter with its energy bound, and an exact
      equivalence between quotient positivity and contractive channel
      factorization.
\item The Meyer route supplies the actual source-native cohomology,
      reciprocal flow, Xi spectral inclusion, and canonical half-density
      isometric completion, including exact graph-closure phase transition
      and a literal M\"obius--Gaussian Neumann source before continuation.
\item Recurrence laws count compact projective directions independently of
      split slope; HT Cartan comparison proves higher-rank invariance under
      finite differential re-embedding.
\end{itemize}

Their common conclusion is structural: downstream augmentation does not
control a resolved source packet unless its actual sector is conservative;
local compatibility does not imply global gluing; and topology does not
choose Hermitian scale.  All finite source, reporter, local incidence,
orientation-detector, face-gluing, and positive-consumer data required
before global descent exist.  Remaining arrow must be arithmetic and
source-pointed: identify actual Green factor provenance with local
incidence, glue it uniquely through the nonlocal singular completion, and
derive the contractive four-unit polarization.  An equivalent source law
may instead force the unsquared trace to vanish.  Only then may the positive
consumer read the zero-Cartan consequence.

\section{Discussion}

\subsection{What the universal law explains}

The parent law is source-pointed polarized exact descent.  Its algebraic
part is the conjunction
\[
       S\cap\ker O=0,
       \qquad O(S)=A\cap\ker\delta:
\]
source lifts are unique and every compatible admissible observation glues.
Its topological part is a positive lower modulus.  Its Hermitian part is
Gram domination, equivalently a contractive factor.  Serial and parallel
network laws propagate the first gate; augmentation is its conservation
specialization; normal-crossing and Sullivan constructions realize the
second; graph verticality and Rees growth measure failure of the third;
Douglas factorization identifies the positive gate.  None implies another
without a source theorem.

Cartan-regular transport is the rank-one character shadow: invariant
positive descent is the zero fibre.  For a nonzero slope the realized
failures are radicalization, reciprocal condition growth, opposite-slope
cancellation or indefinite pairing, singular grading, full verticality,
and factorization failure.  Source entropy, Xi residue, nearby-cycle charge,
Green defect, Borel/Sonin exponent, Rees type, Burnol holonomy, and Meyer
graph closure are coordinate readers or gate tests for one resolved source
packet.  Bounded-prefix Cartan comparison is its higher-rank transport
form; projective recurrence belongs to the independent compact character
coordinate.

The law separates three questions often conflated in RH research.

\begin{enumerate}[leftmargin=2.2em]
\item \emph{Is the divisor represented?}  Yes: several source
      constructions reconstruct it, and Meyer cohomology contains every
      Xi zero with multiplicity.
\item \emph{Is the critical-line orientation represented?}  Yes:
      \(\kappa\) is an explicit projective source character, residue,
      nearby-cycle charge, and resolved Green defect.  A nonzero leading Xi
      jet now carries that displacement into an unsquared completed Hodge
      trace which vanishes exactly at zero displacement.
\item \emph{Why must the orientation vanish?}  Total Green covariance does
      not answer this: it is unconditional antidiagonal cancellation.
      Nor does local Hodge recovery: it works for every displacement.
      Required theorem is the still-missing global conservative edge in
      \eqref{eq:missing-resolved-sector}, or a source-derived vanishing law
      for the unsquared trace after identifying it with Green orientation.
      Its scalar
      shadows are boundedness of the centered evaluation orbit of one
      target-free Gaussian, boundedness of an explicit reflection-fixed
      family of uniformly bounded source capacity, and continuity for the
      canonical half-density quotient seminorm.  None follows from
      hyperbolic conservation, recurrence rank, scalar covariance, local
      positivity, or regular completion.
\end{enumerate}

The new point is that this last question is not one unspecified inequality.
It is failure of one global polarized exact-descent diagram.  Local
existence is witnessed by theta incidence, face gluing, and the Neumann
source; raw completed stability is excluded by nonlocal--singular and
full-verticality theorems; numerical calibration is not encoded by source
topology and is exactly a contractive factorization problem.  This four-stage
audit has a forced order: exact removal of redundant compatibility data
must occur before Hermitian calibration.  It locates both what is available
and what no existing downstream consumer can supply.

This explains why increasingly exact divisor reconstructions did not move
zero exclusion.  They solved representation repeatedly.  Orientation and
all its condition readers are now explicit.  Parent law says no downstream
regular construction can convert those data into purity; purity must enter
as arithmetic support at source.

\subsection{Scope}

All theorems stated as results in this paper are unconditional.  Resolved
arithmetic source support with completed augmentation conservativity---and
hence its zero-Cartan fixed-orbit and half-density consequences---is not
claimed.  The paper does not prove RH, simplicity of zeros, a trace formula,
or global Meyer--Xi spectral synthesis.

Rank-one identities are elementary in isolation.  Their force here is
closure under every realized transport category and their exact arithmetic
realization: cofinal collapse with injective normal tomography, matching
Borel/Sonin exponents, odd hyperbolic holonomy, maximal off-line Xi graph
verticality, exact finite-fibre escape, higher-rank differential Cartan
comparison, independent recurrence rank, and a source-native cohomology
with full finite Xi jet packets.

\section{Conclusion}

Every source-pointed exact descent has three independent obligations:
unique source separation, gluing of every compatible admissible datum, and
stable polarized recovery.  They also have a forced order: on a total test
packet, coefficient preservation fixes the reporter, so strict Hermitian
recalibration is impossible before a nontrivial exact descent.  Serial edges
propagate separation, parallel readers intersect radicals, and Hilbert Gram order is contractive
factorization.  Conservation does not imply absence of internal defect; its
augmentation edge says only that the resolved defect has zero coordinate
sum.  This polarized exact-descent law is the universal law behind the
constructions studied here.  The complete Riemann source
already carries the full complex spectral character.  Its split Cartan
coordinate is conserved as prime-shell entropy, Xi charge, nearby-cycle
charge, Green defect, exact condition exponent, and odd holonomy.  Its
compact coordinate survives separately through projective recurrence.
Finite differential re-embedding preserves the corresponding normalized
higher-rank Cartan vector.

Off axis, positive descent radicalizes the split character and the canonical
Xi graph closure becomes maximally vertical; ordinary completion cannot
repair that failed edge.  Meyer cohomology, a uniform-capacity reflected
Gaussian, one common finite Hardy reporter, coefficient-faithful
physical/Hodge incidence, exact normal-crossing gluing, and a literal
M\"obius--Gaussian Neumann source show that source capacity, local
incidence, algebraic gluing, orientation detection, and finite nodewise
compatibility are already available.  Four-unit Euler capacity is also
available; the division-free completed reporter instead carries eight, and
totality proves that no same-carrier coefficient-preserving replacement can
close the gap.

Exact frontier is ordered polarized singular arithmetic descent.  Canonical
half-density completion exists, and an actual Xi graph is single-valued in
it exactly on the critical line; off line its vertical fibre is all of
\(\C\).  Minimal equivalent is boundedness of one explicit Gaussian
evaluation orbit.  Canonical scalar form is inclusion of every actual Xi
functional in the continuous zero-Cartan support.  Parent law shows what
that scalar form must resolve and in which order: a global source-derived
comparison of actual Green factor defects with coefficient-faithful Hodge
incidence, exact removal of compatibility boundaries, unique gluing through
the forced nonlocal singular completion, and only then a Hermitian Gram
domination producing the correctly normalized contractive factor.
Equivalently, one must derive vanishing of the unsquared secondary
trace from source structure, not assume zero displacement.

\end{document}
